% concatenated from main_arxiv.tex, paper.tex, appendix.tex
\documentclass{article}

% Bundle PRIMEarxiv.sty beside this file in the arXiv source archive.
\usepackage{PRIMEarxiv}

\usepackage[T1]{fontenc}
\usepackage[authoryear,round]{natbib}
\usepackage{url}
\usepackage{booktabs}
\usepackage{array}
\usepackage{tabularx}
\usepackage{graphicx}
\usepackage{adjustbox}
\usepackage{amsmath,amssymb,amsfonts,mathtools}
\usepackage{bm}
\usepackage{amsthm}
\usepackage{microtype}
\usepackage{xcolor}
\usepackage{enumitem}
\usepackage{placeins}
\usepackage{float}
\usepackage{needspace}
\usepackage{algorithm}
\usepackage{algpseudocode}
\usepackage{hyperref}

\hypersetup{
  hypertexnames=false,
  colorlinks=true,
  linkcolor=blue!45!black,
  citecolor=blue!45!black,
  urlcolor=blue!45!black,
  pdftitle={From Truncation to Commitment: Persistent Context in Uniform Discrete Diffusion},
  pdfauthor={Satoshi Hayakawa}
}

\newcommand{\KL}{D_{\mathrm{KL}}}
\newcommand{\kldiv}{\,\Vert\,}
\newcommand{\TV}{\operatorname{TV}}
\newcommand{\TC}{\operatorname{TC}}
\newcommand{\Ent}{\operatorname{Ent}}
\newcommand{\E}{\mathbb E}
\newcommand{\1}{\mathbf 1}
\newcommand{\CRS}{\textsc{CRS}}
\newcommand{\Hinf}{H{=}\infty}

\newtheorem{theorem}{Theorem}
\newtheorem{proposition}[theorem]{Proposition}
\newtheorem{corollary}[theorem]{Corollary}
\newtheorem{lemma}[theorem]{Lemma}
\theoremstyle{definition}
\newtheorem{definition}[theorem]{Definition}
\newtheorem{remark}[theorem]{Remark}

\newcolumntype{Y}{>{\raggedright\arraybackslash}X}

\title{From Truncation to Commitment:\\
Persistent Context in Uniform Discrete Diffusion}

\author{
  Satoshi Hayakawa \\
  The University of Tokyo \\
  \texttt{hayakawa@mist.i.u-tokyo.ac.jp}
}

\begin{document}
\raggedbottom

\maketitle
\begin{abstract}
Uniform-state discrete diffusion models update all tokens in parallel while
keeping every position revisable. Even when the commonly used top-$p$
rule leaves only one candidate at a position, that choice affects
only the current reverse step and can be revised at the next sampling step. We ask
what changes when selected hypotheses instead become persistent context for
later predictions. We therefore propose \emph{committed reveal sampling}
(\CRS), a training-free sampler that stores selected
argmax tokens and inserts them into subsequent model inputs.
Our analysis gives a rationale for selecting later and for keeping selected
tokens visible. Under the exact forward process, the Bayes error of selecting a
clean token cannot increase as noise decreases, while in a simple latent-mode
model, keeping the selected token visible helps
later parallel predictions agree on the same sequence-level choice.
Empirically, paired experiments on Duo-distilled then separate
this persistent effect from single-step top-$p$ restriction and scalar
temperature scaling. Under the same finalization rule, \CRS{} without top-$p$
truncation reaches lower generative perplexity (GenPPL) than fixed $p=0.95$ and
$p=0.9$ baselines across budgets of 8--64 function evaluations (NFE). At 64 NFE, the
comparison at matched unigram entropy also gives lower GenPPL for \CRS{},
yielding a more favorable GenPPL--entropy tradeoff. Base Duo shows the same
direction in a descriptive comparison, while other diversity and continuation
metrics can rank these operating points differently. These results identify support
restriction and persistent context as distinct controls of that tradeoff.
\end{abstract}

\section{Introduction}
\label{sec:intro}

Discrete diffusion models generate a sequence through repeated parallel token
updates~\citep{austin2021structured,campbell2022continuous,lou2024discrete,
sahoo2024simple}. Masked models gradually replace an absorbing mask with visible
tokens, and each reveal supplies context to later updates and fixes part of the
generation order~\citep{chang2022maskgit,hayakawa2026demystifying}. In contrast,
uniform-state models such as Duo use vocabulary symbols as noisy states and
keep every position revisable~\citep{sahoo2025duo}. Such revisability supports
continued sequence-wide revision, but removes the explicit record of earlier
decisions.

In practice, uniform samplers often restrict their predictions over clean
tokens before each reverse update. The released Duo sampler uses top-$p$, or
nucleus, sampling to restrict its clean-token predictions before each reverse
update~\citep{holtzman2020curious,deschenaux2026psi}. Top-$p$ narrows the
clean-token support used by a single update. Late in the reverse process, this
support often becomes a singleton, but even then the selected label is not a
reveal: it is used for that update only, and the next sampling step recomputes its
prediction from a revisable noisy state.

This temporal limitation motivates \emph{committed reveal sampling} (\CRS), which uses selected clean-token hypotheses as persistent context. \CRS{} is a sampling algorithm for frozen model parameters that wraps the positionwise uniform-diffusion denoiser. After a warmup with the native revisable sampler, it selects positions with concentrated predictions, stores their argmax tokens, and writes those tokens into later model inputs; unselected positions continue to follow the native sampler.

The central intuition is that parallel updates can combine individually plausible tokens from incompatible sequence-level alternatives. A reliable visible token can reduce this ambiguity and steer later predictions toward the same alternative, whereas an incorrect selection can steer them toward the wrong one. This raises two questions: when should a token be selected, and how does keeping it visible affect later predictions?

We study these questions theoretically and empirically. Theoretically, under the exact forward process, the optimal clean-token selection error cannot increase as noise decreases, providing a reason for the initial warmup phase; a latent-mode analysis then shows how persistent visibility can coordinate later predictions around the selected alternative. Empirically, paired experiments compare CRS with top-$p$ samplers and measure the resulting tradeoff between generative perplexity (GenPPL) and unigram entropy across compute budgets on the base and distilled Duo checkpoints \citep{sahoo2025duo}.

The paper makes three contributions.
\begin{enumerate}[leftmargin=*,itemsep=2pt,topsep=3pt]
  \item \textbf{Persistence and support restriction.} We distinguish the support restriction induced by top-$p$, which applies only to the current sampling step, from persistent context that remains visible in later steps. This distinction motivates \CRS{}, which uses selected clean-token hypotheses as persistent context to guide later token predictions.
  \item \textbf{Theoretical analysis.} Our analysis provides a rationale for two design choices in \CRS{}: warmup phase before token selection and keeping selected tokens visible afterward. Under the exact forward process, waiting until the state is less noisy cannot increase the optimal clean-token selection error (Proposition~\ref{prop:warmup-blackwell}). In a simple latent-mode model, persistent visible context helps later positionwise predictions favor the same sequence-level alternative (Theorem~\ref{thm:shared-mode}).
  \item \textbf{Experiments.} Under a common finalization rule, \CRS{} without
  top-$p$ truncation has lower GenPPL than fixed global
  $p\in\{0.95,0.9\}$ baselines on Duo-distilled across 8--64 NFE. Base Duo
  shows the same ordering in a descriptive comparison. At 64 NFE,
  matched-entropy comparisons also favor \CRS{} on Duo-distilled.
\end{enumerate}

\section{Related Work}
\label{sec:related-work}

\paragraph{Uniform-state discrete diffusion.}
Duo and its leave-one-out (LOO) analysis define the uniform reverse interface used here
\citep{sahoo2025duo,gourevitch2026uniform}.
\mbox{\citet{deschenaux2026psi}} report Duo top-$p$ sweeps with an ancestral endpoint while developing predictor--corrector samplers. Our study asks a different question: we compare support restriction applied only in the current sampling step with persistent context visible in later steps, using the same coordinatewise argmax decoder in the final step for all methods. Because the endpoint protocols differ, the resulting values are not directly comparable.

\paragraph{Joint distribution modeling.}
The native parallel interface constructs each update from a product of coordinatewise predictions; we refer to this interface as a product denoiser. Prior work enriches this interface through correlation-aware mixtures, energy corrections, coupled transitions, or auxiliary simplex-valued reverse states
\citep{hayakawa2025distillation,xu2024edlm,li2026codd,sakurai2026simplax}.
These approaches modify the learned reverse process, its training procedure, or its state representation to capture dependencies beyond coordinatewise product predictions. \CRS{} instead keeps the checkpoint fixed and changes the context available across sampling steps by exposing selected clean-token hypotheses to later denoiser inputs.

\paragraph{Masked diffusion and commitment.}
Masked diffusion turns an absorbing mask into visible clean tokens, so reveal
order is part of the sampler~\citep{sahoo2024simple,shi2024simplified}.
MaskGIT and later methods select or remask positions using confidence or entropy
\citep{chang2022maskgit,hayakawa2026demystifying,benhamu2025entropy,
wang2025remasking}. Fast-dLLM commits confident tokens for parallel decoding
and cache reuse~\citep{wu2025fastdllm}, while deferred commitment postpones
uncertain blockwise decisions~\citep{shu2026deferred}. Other policies use future self-consistency, trajectory stability, or coordination among dependent token values to determine which token predictions to commit and when
\citep{wang2026selfcontained,sun2026pathmatters,wang2026tacg,yao2026cocommit}.
\CRS{} asks a different question: what happens when a stored clean-token hypothesis remains visible across steps in an otherwise fully revisable uniform-state sampler?

\paragraph{Generalized and larger-scale models.}
The masked-diffusion view also underlies large-scale models such as LLaDA and
Dream~\citep{nie2025llada,ye2025dream}.
GIDD generalizes discrete diffusion to interpolating kernels that can combine
masking and uniform noise~\citep{vonrutte2025gidd}. Sumi scales a pure-uniform
instantiation of this framework with a time-agnostic bidirectional
Transformer~\citep{ye2026sumi}. They provide a complementary testbed to our
explicitly time-indexed Duo model and its LOO prediction
interface. Extending persistent-context controls to that setting is a natural
next step.

\section{Uniform Reverse Sampling and
\texorpdfstring{Top-$p$ Support Restriction}{Top-p Support Restriction}}
\label{sec:uniform-reverse}

We first derive the native reverse interface and locate top-$p$ within one
reverse step. This identifies which clean-token distribution top-$p$ restricts
and why even a singleton choice is not carried explicitly into the next
step. Section~\ref{sec:persistent-context} then introduces persistent context
at that interface.

\subsection{From the LOO law to the native reverse step}

Let $\mathcal S=[V]=\{1,\ldots,V\}$ be the vocabulary and
$[L]=\{1,\ldots,L\}$ be the set of positions. Let $q_0$ be the data distribution on
$\mathcal S^L$, and let $(\bm X_t)_{0\le t\le T}$ be a coordinatewise Markov forward process whose transition kernels belong to the uniform-corruption family. In particular,
$\bm X_0\sim q_0$ and
$\bm X_t\mid\bm X_0\sim q_{t\mid0}(\cdot\mid\bm X_0)$. For uniform
corruption, this conditional law factorizes as
\begin{equation}
 q_{t\mid0}(\bm x_t\mid\bm x_0)
 =\prod_{i=1}^Lq_{t\mid0}^i(x_t^i\mid x_0^i),
 \qquad
 q_{t\mid0}^i(z\mid x)
 =(1-\alpha_t)\1\{z=x\}+\frac{\alpha_t}{V}.
 \label{eq:uniform-forward-main}
\end{equation}
Here $\alpha_t$ is monotonically increasing,
with $\alpha_t=0$ giving clean data and $\alpha_t=1$ giving uniform noise,
and $\bm x=(x^i)_{i=1}^L$. Let $\bm x^{-i}$ denote the sequence with
coordinate $i$ removed.

%To describe a reverse step from $t$ to an earlier time $s$, suppose
For $0\le\alpha_s<\alpha_t\le1$, Markov consistency requires
$q_{t|0}^i(z\mid x) = \sum_y q_{t|s}^i(z\mid y)q_{s|0}^i(y\mid x)$.
Since uniform channels are closed under composition, $q_{t|s}$ is given as
\begin{equation}
 q_{t\mid s}^i(z\mid y)
 =(1-\beta_{t,s})\1\{z=y\}+\frac{\beta_{t,s}}{V},
 \qquad
 \beta_{t,s}=\frac{\alpha_t-\alpha_s}{1-\alpha_s}.
 \label{eq:intermediate-uniform-main}
\end{equation}
The reverse update also needs the earlier noisy token conditioned on both the
clean token and the current noisy token. %Writing $x=x_0^i$ and $y=x_s^i$,
Bayes' rule gives this exact coordinate bridge:
\begin{equation}
 q_{s\mid0,t}^i(y\mid x,z)
 =\frac{q_{s\mid0}^i(y\mid x)q_{t\mid s}^i(z\mid y)}
 {q_{t\mid0}^i(z\mid x)}.
 \label{eq:bridge-main}
\end{equation}

The released Duo model takes the full noisy sequence $\bm x_t$ as input and
outputs, at each coordinate $i$, a normalized categorical law
$h_{\theta,t}^i(\cdot\mid\bm x_t)$ over clean labels
\citep{sahoo2025duo}. Subsequent analysis by
\citet{gourevitch2026uniform} showed that the corresponding population target
is the leave-one-out (LOO) conditional, which predicts the clean token at
coordinate $i$ from all other noisy coordinates:
\begin{equation}
 q_{0\mid t}^i(x\mid\bm x_t^{-i})
 :=\mathbb P(X_0^i=x\mid\bm X_t^{-i}=\bm x_t^{-i})
 =\mathbb E\!\left[
 q_{0\mid t}^i(x\mid\bm X_t)
 \,\middle|\,
 \bm X_t^{-i}=\bm x_t^{-i}
 \right].
 \label{eq:loo-main}
\end{equation}
We therefore interpret $h_{\theta,t}^i$ as a learned approximation to this
population target, while allowing a finite checkpoint to retain some
dependence on the $i$-th coordinate of input $\bm{x}$.

To use the LOO law in the reverse update, we first convert it to the ordinary
denoising posterior. Bayes' rule gives
\begin{equation}
 q_{0\mid t}^i(x\mid\bm x_t)
 =\frac{q_{0\mid t}^i(x\mid\bm x_t^{-i})q_{t\mid0}^i(x_t^i\mid x)}
 {\sum_{\bar x}q_{0\mid t}^i(\bar x\mid\bm x_t^{-i})q_{t\mid0}^i(x_t^i\mid\bar x)}
 =:\bigl(\mathcal C_{t,x_t^i}q_{0\mid t}^i(\cdot\mid\bm x_t^{-i})\bigr)(x),
 \label{eq:conversion-main}
\end{equation}
where, for any categorical law $h$,
$(\mathcal C_{t,z}h)(x):=
h(x)q_{t\mid0}^i(z\mid x)/
\sum_{\bar x}h(\bar x)q_{t\mid0}^i(z\mid\bar x)$.
Indeed, conditional on $X_0^i=x$, corruption at coordinate $i$ is independent
of $\bm X_t^{-i}$ and has law $q_{t\mid0}^i(\cdot\mid x)$.
% Therefore,
% Equation~\eqref{eq:conversion-main} states compactly that
% $q_{0\mid t}^i(\cdot\mid\bm x_t)=
% \mathcal C_{t,x_t^i}q_{0\mid t}^i(\cdot\mid\bm x_t^{-i})$.

Finally, marginalizing the bridge under this ordinary denoising posterior
gives the exact reverse marginal
\begin{equation}
 q_{s\mid t}^i(y\mid\bm x_t)
 =\sum_x q_{s\mid0,t}^i(y\mid x,x_t^i)
 q_{0\mid t}^i(x\mid\bm x_t).
 \label{eq:exact-reverse-main}
\end{equation}
The learned reverse transition $\pi_{s|t}$ substitutes the converted learned law for the
exact posterior:
\begin{equation}
 \pi_{s\mid t}^i(y\mid\bm x_t)
 :=\sum_x q_{s\mid0,t}^i(y\mid x,x_t^i)
 \bigl(\mathcal C_{t,x_t^i}h_{\theta,t}^i
 (\cdot\mid\bm x_t)\bigr)(x),
 \label{eq:learned-reverse-main}
\end{equation}
where the dependence on $\theta$ is omitted from $\pi_{s|t}$ for simplicity.
Equations~\eqref{eq:conversion-main} and~\eqref{eq:learned-reverse-main} separate
three objects: the learned law over clean labels, its local conversion to an
ordinary denoising posterior, and the resulting learned reverse transition,
from which the next noisy token is sampled.

\subsection{\texorpdfstring{Top-$p$}{Top-p} in the LOO distribution}
\label{sec:top-p-loo}

For a categorical law $h$, write
$h(A):=\sum_{a\in A}h(a)$ and, whenever $h(A)>0$,
\begin{equation}
 (\mathsf T_Ah)(a):=\frac{h(a)\1\{a\in A\}}{h(A)}.
 \label{eq:general-restriction}
\end{equation}
Sort the vocabulary as
$h(a_1)\ge\cdots\ge h(a_V)$ with a fixed tie rule. The top-$p$ support
induced by $h$ is
\begin{equation}
  \mathsf T_ph:=\mathsf T_{A_p(h)}h, \quad
  \text{where}\ 
 A_p(h)=\{a_1,\ldots,a_{k_p(h)}\},
 \ 
 k_p(h)=\min\Bigl\{k:\sum_{j=1}^k h(a_j)\ge p\Bigr\}.
 \label{eq:top-p-support-main}
\end{equation}
Top-$p$ acts within one reverse step by substituting $\mathsf T_p h$ for $h$ in
Equation~\eqref{eq:learned-reverse-main} \citep{holtzman2020curious,hewitt2022desmoothing}; the resulting transition samples the
next noisy state (Figure~\ref{fig:concept}(a)). Even when $A_p(h)$ is a
singleton, its label is absent from the next denoiser input unless the noisy
state happens to preserve it.
Empirically, at 64 NFE on Duo-distilled with $p=0.9$, the top-$p$ support is a singleton on
$27.5\%$ of active coordinate--step pairs overall and $45.6\%$ in the final ten
reverse steps (Appendix~\ref{app:singleton-trace}).

The conversion in Equation~\eqref{eq:conversion-main} preserves hard support
restriction. For every retained set $A$ with positive normalizers,
$%\begin{equation}
 \mathcal C_{t,x_t^i}(\mathsf T_Ah)
 =\bigl(\mathcal C_{t,x_t^i}h\bigr)(\cdot\mid A)
$ %\end{equation}
holds.
Indeed, writing $Q=\mathcal C_{t,x_t^i}h$, direct substitution gives
$\mathcal C_{t,x_t^i}(\mathsf T_Ah)(a)=Q(a)\1\{a\in A\}/Q(A)$, which is
$Q(\cdot\mid A)$. Thus top-$p$ is hard support restriction at the clean-label
interface used by the reverse bridge.
In the singleton limit it makes a point decision for the current reverse step, while the
state update can remain stochastic and the point label is recomputed at the
next step.

\begin{figure}[H]
  \centering
  \includegraphics[width=0.98\linewidth]{figures/truncation_commitment_concept.pdf}
  \caption{Support restriction and persistent context act at different points.
  Top-$p$ restricts the clean-token distribution used by one reverse update.
  \CRS{} stores selected clean tokens (squares) and inserts them into later
  denoiser inputs. Other positions remain revisable. Shape denotes whether a
  position remains revisable; colors illustrate current token values.}
\label{fig:concept}
\end{figure}

\section{From Point Hypotheses to Persistent Context}
\label{sec:persistent-context}

Top-$p$ can collapse the clean-token support in one reverse step to a point,
but does not retain that hypothesis as explicit context: the next denoiser
receives only the updated noisy state. A persistent sampler instead keeps
selected clean-token values explicit in later denoiser inputs. \CRS{} does so
while leaving all unselected positions revisable, paralleling the use of
revealed tokens as later context in masked diffusion
\citep{chang2022maskgit,hayakawa2026demystifying}.

\subsection{Committed reveal sampling}
\label{sec:crs-method}

Let $C\subset[L]$ be the selected coordinates and
$\bm c=(c_i)_{i\in C}$ their stored labels. We write
$\bm x[C\leftarrow\bm c]$ for replacing those coordinates in a model input.
During an initial warmup phase, \CRS{} runs the native sampler without storing
labels. It then distributes a total confidence-deficit budget $D_{\max}$ over
the remaining steps. Selecting coordinate $i$ incurs the deficit
$1-\max_a h_{\theta,t}^i(a\mid\widetilde{\bm x}_t)$.
After the warmup phase, each sampling step of \CRS{} is as follows:
\begin{enumerate}[leftmargin=*,itemsep=1pt,topsep=2pt]
  \item evaluate the denoiser on
  $\widetilde{\bm x}_t:=\bm x_t[C\leftarrow\bm c]$, so labels stored in earlier
  steps are visible;
  \item use this output for the native reverse transition, scan the currently
  unselected coordinates from highest to lowest confidence, and greedily
  select a spatially separated batch within the budget allocated to that
  step; and
  \item store the argmax label at each newly selected coordinate and update
  $(C,\bm c)$, making the selected labels visible in subsequent denoiser
  evaluations.
\end{enumerate}

We form each batch using a fixed local-exclusion rule;
Appendix~\ref{app:crs-algorithm} gives the exact constraint and connects it to
spatially dispersed scheduling
\citep{besnier2025halton,hayakawa2026demystifying}.
For each selected coordinate, \CRS{} stores the argmax of its clean-token
prediction. This value coincides with the unique retained token whenever
top-$p$ produces a singleton support. More generally, argmax is the
deterministic choice that minimizes coordinatewise Hamming risk under the
denoiser output:
\begin{equation}
 \arg\min_{c\in\mathcal S}
 \mathbb E_{X\sim h}\!\left[\1\{X\ne c\}\right]
 =\arg\max_{c\in\mathcal S}h(c).
 \label{eq:argmax-hamming-main}
\end{equation}
Sampling from the top-$p$-restricted distribution when its support is not a
singleton, and then storing that draw, would combine token-sampling randomness
with contextual feedback. \CRS{} therefore uses argmax to isolate the latter.
Indeed, at $p=0.9$ and 64 NFE it matches top-$p$ on the $45.6\%$ of active
coordinate--step pairs with singleton support in the final ten reverse steps, and
Equation~\eqref{eq:argmax-hamming-main} gives its coordinatewise Hamming-risk
justification otherwise. The full selector and pseudocode, together with
experiments that limit how long stored labels remain visible, appear in
Appendix~\ref{app:crs-algorithm}.

\paragraph{Common final-step decoder.}
The released Duo sampler uses an ancestral endpoint
\citep{sahoo2025duo,deschenaux2026psi}. To compare support restriction and
persistent context under a shared output rule, our experiments instead count a
coordinatewise argmax readout as the final one of $N$ sampling steps. Each
step uses one denoiser evaluation, so NFE is $N$. Let
$\bm x_{\varepsilon}$ be the last noisy state on
the grid $t_0>\cdots>t_{N-1}=\varepsilon>0$. The first $N-1$ steps produce
stochastic reverse transitions, and the final step returns
$\widehat x^i=\arg\max_{a\in\mathcal S}
h_{\theta,\varepsilon}^i(a\mid\widetilde{\bm x}_{\varepsilon})$ at every
generated coordinate. This comparison protocol uses the model's clean-token
output directly rather than sampling another reverse transition.

For $A\subseteq C$, let $\bm c|_A=(c_i)_{i\in A}$ denote the corresponding
restriction of the stored labels. For CRS,
$\widetilde{\bm x}_{\varepsilon}=\bm x_{\varepsilon}
[C_{\mathrm{vis}}\leftarrow\bm c|_{C_{\mathrm{vis}}}]$, where
$C_{\mathrm{vis}}\subseteq C$ indexes coordinates
with stored labels still visible. The selected tokens appear in this final
model input but are never copied directly to the output. We restore fixed
conditioning positions separately.

\subsection{Why persistent context can coordinate parallel updates}
\label{sec:theory}

Parallel product updates can combine individually plausible tokens from
incompatible sequence-level alternatives. Correlation-aware denoisers and
joint transitions address this by modifying dependence within a step
\citep{hayakawa2025distillation,xu2024edlm,li2026codd}. \CRS{} instead retains
a product output and changes later predictions through visible context.

% \paragraph{Two-mode example.}
To isolate this effect, consider two binary coordinates with clean law
$Q_\eta=(1-\eta)\delta_{00}+\eta\delta_{11}$ for $0<\eta<1/2$.
The product of its marginals assigns total mass $2\eta(1-\eta)$ to the hybrid
strings $01$ and $10$, while $Q_\eta$ assigns them zero mass. For
$X=(X^1,X^2)\sim Q_\eta$, coordinatewise argmax selects $00$, with
sequence-level error $\mathbb P(X\ne00)=\eta$. If a selected value
$c_1\in\{0,1\}$ remains visible at the first coordinate, then
$Q_\eta(X^2=c_1\mid X^1=c_1)=1$; the other coordinate is symmetric.
Persistence, rather than argmax alone, carries this mode evidence into later
predictions.

We now generalize this mechanism. Let a latent mode
$M\in[K]$ have prior weights $w_m=\mathbb P(M=m)$. Conditional on $M=m$,
let the future coordinates $Y=(Y_1,\ldots,Y_n)$ be independent, with laws
$P_{m,i}$. Define
\begin{equation}
 Q_{\mathrm{mix}}^w(y)=\sum_{m=1}^K w_m\prod_{i=1}^nP_{m,i}(y_i),
 \qquad
 Q_{\mathrm{prod}}^w(y)=\prod_{i=1}^n\sum_{m=1}^K w_mP_{m,i}(y_i).
 \label{eq:shared-mode-laws}
\end{equation}
The mixture uses one shared mode, whereas the product law independently
remixes its coordinate marginals. Here $\KL$ denotes Kullback--Leibler (KL)
divergence, $\TC$ denotes total correlation~\citep{watanabe1960information}, a
standard measure of residual multivariate dependence whose conditional variants
have also been used in discrete diffusion~\citep{yoo2025redi}, and
$\Ent(w):=-\sum_m w_m\log w_m$. We also write
$h_2(a):=-a\log a-(1-a)\log(1-a)$.

\begin{theorem}[Shared-mode coordination]
\label{thm:shared-mode}
For the laws in Equation~\eqref{eq:shared-mode-laws},
\begin{equation}
 \KL(Q_{\mathrm{mix}}^w\kldiv Q_{\mathrm{prod}}^w)
 =\TC_{Q_{\mathrm{mix}}^w}(Y_1,\ldots,Y_n)
 \le(n-1)\Ent(w).
 \label{eq:shared-mode-gap}
\end{equation}
Let $U$ be visible context with $Y\perp\!\!\!\perp U\mid M$. For any
positive-probability value $u$, let $w_m[u]=\mathbb P(M=m\mid U=u)$ and
$\epsilon(u):=1-\max_m w_m[u]$. Then
\begin{equation}
 \KL(Q_{\mathrm{mix}}^{w[u]}\kldiv Q_{\mathrm{prod}}^{w[u]})
 \le(n-1)\{h_2(\epsilon(u))+\epsilon(u)\log(K-1)\}.
 \label{eq:context-mode-gap}
\end{equation}
\end{theorem}

Appendix~\ref{app:shared-mode-theorem-proof} gives the proof.
A mode-consistent visible label concentrates $w[u]$ and drives the bound
toward zero; a poor label can instead favor the wrong mode. The appendix gives
an explicit likelihood-ratio condition and extends the calculation to combined
support restriction and persistent evidence
(Appendix~\ref{app:combined-evidence}).

\subsection{Why wait before selecting tokens?}

The CRS warmup introduced in Section~\ref{sec:crs-method} changes two
quantities: it can make the selected value more reliable,
and it leaves fewer later steps in which that value can affect the path, a
tension also reflected in recent confidence- and trajectory-based commitment
policies~\citep{shu2026deferred,sun2026pathmatters}. The
following proposition isolates the first effect at an exact forward state. For
coordinate $i$, define the oracle LOO Hamming Bayes risk at noise level $t$ by
\begin{equation}
 \mathcal R_t^i
 :=\mathbb E\!\left[1-\max_{a\in\mathcal S}
 q_{0\mid t}^i(a\mid\bm X_t^{-i})\right].
 \label{eq:loo-warmup-risk}
\end{equation}

\begin{proposition}[Lower noise weakly reduces oracle selection error]
\label{prop:warmup-blackwell}
For the coordinate-factorized uniform forward process and noise levels
$0\le\alpha_s<\alpha_t\le1$,
$\mathcal R_s^i\le\mathcal R_t^i$ for every coordinate $i$.
\end{proposition}

Appendix~\ref{app:warmup-proposition-proof} gives the proof.
Additional corruption maps $\bm X_s^{-i}$ to $\bm X_t^{-i}$, so the cleaner
context can simulate any predictor based on the noisier one. The proposition
concerns the exact posterior; in practice, we rank positions using the learned
model's confidence deficit. Changing the warmup also changes how many later
steps can use a selected token, so the comparison in Figure~\ref{fig:warmup-permanent} measures both effects together.

\section{Experiments}
\label{sec:experiments}

We use the public Duo-distilled and base Duo checkpoints~\citep{sahoo2025duo}
to study two questions. First, how do top-$p$ and \CRS{} trade generative
quality for entropy across compute budgets? Second, what do paired controls
reveal about the effect of persistent context on later predictions? All
end-to-end comparisons use argmax decoding in the final step, and each policy
generates its own trajectory. Following prior discrete-diffusion
studies~\citep{hayakawa2025distillation,hayakawa2026demystifying}, our main
quality metric is
$\mathrm{GenPPL}=\exp(\mathrm{GenNLL})$, where GenNLL is the mean next-token
negative log-likelihood assigned to generated samples by
GPT-2-Large~\citep{radford2019language}. In the main text, we use mean
within-sequence unigram entropy as the diversity axis.
Appendix~\ref{app:experiment-details} reports complementary diversity and
repetition metrics, full configurations, experiments that limit how long
selected labels remain visible, and all uncertainties.

\subsection{How does the tradeoff change with NFE?}
\label{sec:nfe-scaling}

Panels (a)--(b) of Figure~\ref{fig:nfe-scaling} plot GenPPL against NFE while
holding four procedures fixed. On Duo-distilled, \CRS{} with top-$p$ disabled ($p=1.0$) has lower GenPPL throughout the tested NFE range than native sampling with $p\in\{1.0,0.95,0.9\}$. Base Duo shows the same ordering of point estimates in a descriptive, protocol-aligned comparison.
Appendix~\ref{app:nfe-scaling} gives estimates and uncertainties.

\begin{figure}[H]
  \centering
  \includegraphics[width=\linewidth]{figures/nfe_scaling.pdf}
  \caption{
    GenPPL across NFE (a--b) and 64-NFE GenPPL--entropy comparisons
    (c--d), all with argmax decoding in the final step. Parentheses in
    (a--b) give unigram entropy. Panels (c--d) sweep global top-$p$ and the
    \CRS{} selection budget; open diamonds in (c) show the
    $D_{\max}=512$ warmup path used in Figure~\ref{fig:warmup-permanent}.
    Appendix~\ref{app:nfe-scaling} gives the complete grids and uncertainties.
  }
  \label{fig:nfe-scaling}
\end{figure}

Note that our base-Duo top-$p$ values differ from those reported by
\mbox{\citet{deschenaux2026psi}} because their final step samples ancestrally,
whereas ours uses argmax decoding for every method. In a paired 32-NFE audit
that fixes all earlier reverse steps, changing only the final step back to
ancestral sampling recovers their reported range for $p\in\{1.0,0.9\}$.
Figure~\ref{fig:nfe-scaling} therefore compares the methods under our common
argmax final step. Appendix~\ref{app:nfe-scaling} gives the paired values.

\subsection{What changes when selected tokens remain visible?}
\label{sec:mechanism-evidence}

Three controls show that the path change depends on the inserted token values
and cannot be reduced to position selection or scalar temperature. Replacing
stored argmax labels by nearby alternatives reverses the likelihood gain, so
selected positions alone are insufficient. At the same mutable state, the
best scalar-temperature fit leaves
$79.4\pm2.1\%$ of the pooled KL shift unexplained, and the top-ranked token
changes on $30.2\pm0.8\%$ of audited coordinates. Restoring one
selected coordinate to its revisable input leaves the surrounding state
predictive of that label, indicating that its information has propagated to
other coordinates. Appendix~\ref{app:mechanism-tables} gives
these estimates. Appendices~\ref{app:storage-output-control} and
\ref{app:support-duration} test direct output copying and vary how many later
steps can see each selected label.

\paragraph{How do the checkpoints differ?}
At each shared uniform-forward state, we apply the same \CRS{} position-selection rule and budget separately to the two checkpoints. The resulting selected-coordinate sets overlap only moderately. Conditional on a fixed set of selected coordinates, however, the checkpoints almost always predict the same argmax labels and attain nearly identical clean-label accuracy. Duo-distilled mainly differs by selecting positions with slightly lower oracle error. This shared-state audit therefore localizes a checkpoint difference in \CRS{} position selection. Appendix~\ref{app:checkpoint-selector} gives the factorial and calibration details.

\paragraph{What changes when selection starts later?}
On exact uniform-forward states, selected-label error falls steadily as noise
decreases, consistent with Proposition~\ref{prop:warmup-blackwell}. The
sampler-level sweep varies how many steps precede selection. Later selection
recovers unigram entropy while retaining a substantial likelihood gain, and
all four warmup endpoints lie below the matched global top-$p$ curve in
Figure~\ref{fig:nfe-scaling}(c). Appendix
\ref{app:warmup-duration} gives uncertainties, controls that vary how many
later steps see each selected label, and sequence-level metrics.

\begin{figure}[H]
  \centering
  \includegraphics[width=0.68\linewidth]{figures/warmup_permanent_tradeoff.pdf}
  \caption{The \CRS{} tradeoff as selection starts later at 64 NFE.
  Ticks give steps before first selection; parentheses give GenPPL.
  Appendix~\ref{app:warmup-duration} gives uncertainties and controls that
  vary how many later steps see each selected label.}
  \label{fig:warmup-permanent}
\end{figure}

\subsection{Are support restriction and persistent context interchangeable?}
\label{sec:complete-policy-evidence}

Panels (c--d) of Figure~\ref{fig:nfe-scaling} compare CRS budget
sweeps with global top-$p$ at 64 NFE. Across the overlapping entropy region,
the CRS curve lies below the global curve on Duo-distilled; base Duo shows the
same direction in a descriptive comparison. After matching entropy within each
shard, \CRS{} at $D_{\max}=512$ lowers GenNLL by $0.154\pm0.031$ nats per
token.
Appendices~\ref{app:nfe-scaling} and~\ref{app:mechanism-tables} report the
frontier values, uncertainties, and interpolation rule. Global top-$p$ has
lower repeated 4-gram fraction (Rep-4)~\citep{welleck2020unlikelihood}; in
fixed-prefix continuation, GPT-2 favors applying $p=0.8$ at every step, whereas
Pythia~\citep{biderman2023pythia} does not distinguish the two within
uncertainty.

The curves overlap in places, but the samplers reach them differently:
top-$p$ truncates the distribution used at each current step, whereas \CRS{}
carries selected tokens into later model inputs. Delaying selection then moves
\CRS{} along its GenPPL--entropy curve. Rep-4 and fixed-prefix likelihood rank
some of these operating points differently.

\section{Concluding Remarks}
\label{sec:conclusion}

\CRS{} turns selected clean-token hypotheses into context for later steps,
whereas top-$p$ restricts only the distribution used at the current step. Our
analysis gives the two stages different roles: waiting weakly improves oracle
token selection, while keeping a selected token visible can steer later
product predictions toward the same mode.

Across 8--64 NFE under the same argmax final step, \CRS{} with top-$p$
disabled ($p=1.0$) achieves lower GenPPL than the fixed
$p\in\{0.95,0.9\}$ baselines on both checkpoints. At 64 NFE, varying its
selection budget gives lower GenPPL at matched entropy on Duo-distilled; a
separate base-Duo comparison shows the same trend. This advantage is specific
to the GenPPL--entropy comparison: Rep-4 and fixed-prefix likelihood can rank
the operating points differently. Future work should test larger checkpoints
and other uniform-state architectures.


\section*{Acknowledgments}
The author thanks Chunsan Hong for helpful comments and pointers to relevant literature.
The author received access to ChatGPT and Codex through OpenAI's
ChatGPT for Academic Researchers program.

\clearpage
\bibliographystyle{plainnat}
\bibliography{references}

\clearpage
\appendix
\iffalse
\section{Restriction Identities and Path Ledger}
\label{app:path-ledger}

\subsection{Proof of Proposition~\ref{prop:converted-restriction}}

For any $A\subseteq\mathcal S$ with $h(A):=\sum_{a\in A}h(a)>0$, define
\[
 (\mathsf T_Ah)(a):=\frac{h(a)\1\{a\in A\}}{h(A)}.
\]
Assume also that the two conversion normalizers below are positive. Let
$Z=\sum_a h(a)q_{t\mid0}^i(x_t^i\mid a)$ and
$Z_A=\sum_{a\in A}h(a)q_{t\mid0}^i(x_t^i\mid a)$. Direct substitution gives
\begin{equation}
 Q_A(a)=\frac{h(a)q_{t\mid0}^i(x_t^i\mid a)\1\{a\in A\}}
 {Z_A}
 =\frac{Q(a)\1\{a\in A\}}{Q(A)}.
\end{equation}
Hence $Q_A=Q(\cdot\mid A)$. The total-variation identity follows by summing
the positive and negative parts over $A$ and $A^c$. Moreover,
\begin{equation}
 \KL(Q_A\kldiv Q)
 =\sum_{a\in A}Q_A(a)\log\frac{1}{Q(A)}=-\log\widetilde m,
\end{equation}
and the Bhattacharyya coefficient is
$\sum_a\sqrt{Q_A(a)Q(a)}=\sqrt{\widetilde m}$. This proves all three
identities under the convention
$d_{\mathrm{Hel}}^2(P,Q):=\frac12\sum_a(\sqrt{P(a)}-\sqrt{Q(a)})^2
=1-\sum_a\sqrt{P(a)Q(a)}$.

\subsection{Exact charge and restriction geometry}
\label{app:charge-geometry}

The one-coordinate identity extends exactly to the product clean-label law
used by a parallel reverse step. Let
\[
 Q=\bigotimes_{i=1}^n Q_i,
 \qquad
 A=\prod_{i=1}^n A_i,
 \qquad
 Q^A=\bigotimes_{i=1}^nQ_i(\cdot\mid A_i),
\]
where $Q_i(A_i)>0$, and define
\[
 \chi_i:=-\log Q_i(A_i),
 \qquad B:=\sum_{i=1}^n\chi_i.
\]

\begin{proposition}[Exact rectangular charge]
\label{prop:rectangular-charge}
The restricted product law satisfies
\begin{equation}
 Q^A=Q(\cdot\mid A),
 \qquad Q(A)=e^{-B},
 \label{eq:rectangular-mass-app}
\end{equation}
and therefore
\begin{equation}
 \KL(Q^A\kldiv Q)=B,
 \quad
 \TV(Q^A,Q)=1-e^{-B},
 \quad
 d_{\mathrm{Hel}}^2(Q^A,Q)=1-e^{-B/2}.
 \label{eq:rectangular-identities-app}
\end{equation}
\end{proposition}

\begin{proof}
Product structure gives
$Q(A)=\prod_iQ_i(A_i)=\exp(-\sum_i\chi_i)=e^{-B}$. For
$\bm x\in A$,
\[
 Q^A(\bm x)
 =\prod_i\frac{Q_i(x_i)}{Q_i(A_i)}
 =\frac{Q(\bm x)}{Q(A)},
\]
and $Q^A$ vanishes outside $A$. Thus $Q^A=Q(\cdot\mid A)$, and
Proposition~\ref{prop:converted-restriction} gives
Equation~\eqref{eq:rectangular-identities-app}.
\end{proof}

Matching $B$ therefore matches retained product mass and the three displayed
interface distances. It does not match which coordinates are restricted or the
shape of the retained rectangle.

\subsection{Entropy disturbance at matched charge}
\label{app:entropy-geometry}

For $m=Q(A)\in(0,1)$, the entropy chain rule gives the exact decomposition
\begin{equation}
 \Ent(Q)=h_2(m)+m\Ent(Q\mid A)+(1-m)\Ent(Q\mid A^c),
 \label{eq:entropy-decomposition-app}
\end{equation}
and hence
\begin{equation}
 \Ent(Q)-\Ent(Q\mid A)
 =h_2(m)+(1-m)\{\Ent(Q\mid A^c)-\Ent(Q\mid A)\}.
 \label{eq:entropy-difference-app}
\end{equation}
The retained set is a descending prefix of the LOO-coordinate denoiser output, while the
entropy above belongs to the converted clean law. Conversion can change token
order, so the signed difference in Equation~\eqref{eq:entropy-difference-app}
can have either sign.

For a two-sided certificate on an alphabet of size $V\ge2$, define
\begin{equation}
 \Phi_V(\chi):=
 \begin{cases}
 h_2(1-e^{-\chi})+(1-e^{-\chi})\log(V-1),
 &0\le\chi\le\log V,\\[1mm]
 \log V,&\chi>\log V.
 \end{cases}
 \label{eq:entropy-envelope-app}
\end{equation}

\begin{proposition}[Entropy disturbance from restriction]
\label{prop:entropy-disturbance}
For every categorical law $Q$ and retained set $A$ with $Q(A)>0$,
\begin{equation}
 \left|\Ent(Q(\cdot\mid A))-\Ent(Q)\right|
 \le \Phi_V(-\log Q(A)).
 \label{eq:entropy-disturbance-app}
\end{equation}
The function $\Phi_V$ is nondecreasing and concave on $[0,\infty)$.
\end{proposition}

\begin{proof}
Proposition~\ref{prop:converted-restriction} gives total variation
$\delta=1-e^{-\chi}$. For $\delta\le1-1/V$, the sharp
Fannes--Audenaert inequality gives
$|\Ent(P)-\Ent(Q)|\le h_2(\delta)+\delta\log(V-1)$
\citep{audenaert2007sharp}. This condition is equivalent to
$\chi\le\log V$. For larger $\delta$, use the universal bound $\log V$.

On $[0,1-1/V]$, the Fannes--Audenaert right-hand side is concave and
nondecreasing in $\delta$. The map $1-e^{-\chi}$ is also concave and
nondecreasing. Their composition is therefore concave and nondecreasing until
$\chi=\log V$, where its derivative is zero. The constant continuation
preserves both properties.
\end{proof}

\begin{corollary}[Concentrated charge and worst-case entropy disturbance]
\label{cor:charge-majorization}
Let $\boldsymbol\chi,\boldsymbol\chi'\in\mathbb R_+^n$ have equal sums. If
$\boldsymbol\chi\succ\boldsymbol\chi'$ in the majorization order, then
\begin{equation}
 \sum_i\Phi_V(\chi_i)
 \le
 \sum_i\Phi_V(\chi_i').
 \label{eq:charge-majorization-app}
\end{equation}
\end{corollary}

\begin{proof}
Apply Karamata's inequality to the concave function $\Phi_V$.
\end{proof}

This result compares a worst-case sum of marginal clean-interface entropy
disturbances. It does not determine their signs or the entropy of generated
text.

\subsection{Bridge contraction and accumulated charge}

Assume $V\ge2$ and $0\le\alpha_s<\alpha_t\le1$. Write
$\beta_{t,s}:=(\alpha_t-\alpha_s)/(1-\alpha_s)$ and, for fixed current token
$x_t^i$, regard
\[
 K_{t,s}^{x_t^i}(y\mid x)
 :=q_{s\mid0,t}^i(y\mid x,x_t^i)
\]
as a channel from the clean token $x$ to the next noisy token $y$. For a finite
channel $K$, its largest input-independent common mass is
\begin{equation}
 \kappa(K):=\sum_y\min_xK(y\mid x).
 \label{eq:common-mass-app}
\end{equation}
Put $a_\gamma:=1-\gamma+\gamma/V$ and $b_\gamma:=\gamma/V$. Directly from
Equation~\eqref{eq:bridge-main},
\begin{equation}
 \kappa(K_{t,s}^{x_t^i})
 =\frac{b_{\alpha_s}a_{\beta_{t,s}}}{b_{\alpha_t}}
 +(V-1)\frac{b_{\alpha_s}b_{\beta_{t,s}}}{a_{\alpha_t}}.
 \label{eq:uniform-common-mass-app}
\end{equation}
Substitution of $\beta_{t,s}$ shows that
$1-\kappa(K_{t,s}^{x_t^i})=\rho_{t,s}$, where
\begin{equation}
 \rho_{t,s}
 =\frac{(\alpha_t-\alpha_s)
 [V-(V-1)\alpha_t-\alpha_s]}
 {(1-\alpha_s)\alpha_t[V-(V-1)\alpha_t]}\in[0,1].
 \label{eq:rho-app}
\end{equation}
This quantity is independent of $x_t^i$. Writing
$\kappa=\kappa(K_{t,s}^{x_t^i})$, for $0<\kappa<1$ define
\begin{align}
 \eta_{t,s}^{x_t^i}(y)
 &:=\frac{\min_xK_{t,s}^{x_t^i}(y\mid x)}{\kappa},\nonumber\\
 \widetilde K_{t,s}^{x_t^i}(y\mid x)
 &:=\frac{K_{t,s}^{x_t^i}(y\mid x)
 -(1-\rho_{t,s})\eta_{t,s}^{x_t^i}(y)}{\rho_{t,s}}.
 \label{eq:uniform-bridge-mixture-parts-app}
\end{align}
The second numerator is nonnegative and sums to $\rho_{t,s}$, hence
\begin{equation}
 K_{t,s}^{x_t^i}(\cdot\mid x)
 =(1-\rho_{t,s})\eta_{t,s}^{x_t^i}(\cdot)
 +\rho_{t,s}\widetilde K_{t,s}^{x_t^i}(\cdot\mid x).
 \label{eq:uniform-bridge-mixture-app}
\end{equation}
The null-step and clean-endpoint cases follow by continuity
($\rho=0$ and $\rho=1$, respectively).

For a clean-label law $R$, write
$(KR)(y):=\sum_xK(y\mid x)R(x)$. Joint convexity and then data processing give
\[
 \KL(KR\kldiv KS)
 \le \rho_{t,s}\KL(\widetilde K R\kldiv\widetilde K S)
 \le \rho_{t,s}\KL(R\kldiv S).
\]
Taking $R=Q_A$ and $S=Q$ therefore yields
\begin{equation}
 \KL(KQ_A\kldiv KQ)\le \rho_{t,s}\KL(Q_A\kldiv Q)
 =\rho_{t,s}(-\log\widetilde m).
 \label{eq:bridge-contraction-app}
\end{equation}

Consider two path laws with the same initial distribution and overwrite rule,
differing in top-$p$ versus full clean-token support. Let
$\mathbb P_p^{\mathrm{path}}$ and $\mathbb P_1^{\mathrm{path}}$ denote their
path laws, $\mathcal H_r$ the history at step $r$, and $I_r$ the updated
coordinates. Their conditional path kernels obey the exact chain-rule identity
\begin{equation}
 \KL(\mathbb P_p^{\mathrm{path}}\kldiv\mathbb P_1^{\mathrm{path}})
 =
 \E_{\mathbb P_p^{\mathrm{path}}}\!\left[
 \sum_r
 \KL\!\left(
 \mathbb P_p(\bm X_{r+1}\mid\mathcal H_r)
 \kldiv
 \mathbb P_1(\bm X_{r+1}\mid\mathcal H_r)
 \right)\right].
 \label{eq:path-chain-rule-app}
\end{equation}
Conditional on a history, the released parallel update is a product over
$i\in I_r$; unchanged coordinates have the same deterministic kernel under
both laws. Product additivity followed by
Equation~\eqref{eq:bridge-contraction-app} therefore gives
\begin{equation}
 \KL(\mathbb P_p^{\mathrm{path}}\kldiv\mathbb P_1^{\mathrm{path}})
 \le
 \E_{\mathbb P_p^{\mathrm{path}}}\!\left[
 \sum_r\sum_{i\in I_r}\rho_r
 (-\log\widetilde m_{r,i})\right].
 \label{eq:path-ledger-app}
\end{equation}
Here $\rho_r:=\rho_{t_r,t_{r+1}}$. Data processing gives the same upper bound
for endpoint KL when both laws use the same deterministic terminal denoiser
decoding. Every summand can be computed from the denoiser output, retained set,
and schedule. In the audited
64-NFE native $p=0.8$ diagnostic, top-$p=0.8$ is applied throughout a
no-overwrite path and compared with full support. Each categorical
conditional-kernel KL in
Equation~\eqref{eq:path-chain-rule-app} is summed exactly at the visited
history; only the outer path expectation is estimated by Monte Carlo. The
result is $603.0\pm20.8$ nats, versus an upper ledger of
$1507.7\pm22.6$ nats per path (mean$\pm$SE over four seed shards, four paths
per shard). We use the ledger as a scale and allocation audit.
\fi

\FloatBarrier

\section{Committed Reveal and Finite Persistence}
\label{app:crs-algorithm}

Let $N$ be the total NFE and let
$t_0>\cdots>t_{N-1}=\varepsilon>0$ be the noisy-state grid. The first $N-1$
sampling steps produce reverse transitions, and the final step is the common
argmax decoder from Section~\ref{sec:persistent-context}. Let $W$ be the
number of warmup steps among the first $N-1$. For a requested warmup fraction
$f$, the implementation sets
$W=\operatorname{round}\{f(N-1)\}$ and uses $R=N-1-W$ post-warmup steps.
At reverse step $r$,
$\bm x_{t_r}$ denotes the mutable noisy state and
$C_r$ the coordinates whose stored labels are currently visible. Our implementation
uses an \emph{input-overwrite-only} interface: it overwrites only the denoiser input,
\[
 \widetilde{\bm x}_r=\bm x_{t_r}[C_r\leftarrow\bm c|_{C_r}].
\]
It does not pin the noisy-state token used by the reverse transition. Thus, after any scheduled
top-$p_r$ restriction, the actual one-coordinate transition is
\begin{equation}
 \pi_{t_{r+1}\mid t_r}^{i,\mathrm{ow}}
 (y\mid\widetilde{\bm x}_r;x_{t_r}^i)
 :=
 \sum_x q_{t_{r+1}\mid0,t_r}^i(y\mid x,x_{t_r}^i)
 \Bigl(
 \mathcal C_{t_r,x_{t_r}^i}
 \mathsf T_{p_r}h_{\theta,t_r}^i
 (\cdot\mid\widetilde{\bm x}_r)
 \Bigr)(x).
 \label{eq:overwrite-transition-app}
\end{equation}
The superscript ``ow'' denotes this input-overwrite interface. The semicolon
separates the overwritten denoiser input from the mutable current-token
argument: $h_{\theta,t_r}^i$ sees $\widetilde{\bm x}_r$, whereas the conversion
and bridge in Equations~\eqref{eq:conversion-main}
and~\eqref{eq:learned-reverse-main} continue to use $x_{t_r}^i$. Here
$\mathsf T_{p_r}h:=\mathsf T_{A_{p_r}(h)}h$ and $\mathsf T_1$ is the
identity. For each eligible coordinate,
\begin{equation}
 a_r^i\in\arg\max_{a\in\mathcal S}h_{\theta,t_r}^i(a\mid\widetilde{\bm x}_r),
 \qquad
 d_r^i=1-h_{\theta,t_r}^i(a_r^i\mid\widetilde{\bm x}_r).
 \label{eq:selector-deficit-app}
\end{equation}
The canonical selector scans currently unselected coordinates in increasing
$d_r^i$ and greedily forms a within-step batch $S_r$. With fixed spacing
radius $R_{\mathrm{sp}}$, every accepted batch satisfies
\begin{equation}
 |S_r|\le K_r,
 \qquad
 \sum_{i\in S_r}d_r^i\le E_r,
 \qquad
 |i-j|>R_{\mathrm{sp}}
 \quad\text{for all distinct }i,j\in S_r.
 \label{eq:selector-constraints-app}
\end{equation}
The spacing test is applied among coordinates newly selected at step $r$;
coordinates stored at earlier steps are excluded from reselection but do not
exclude their neighboring positions at later steps. This local-exclusion rule
is a simple one-dimensional analogue of spatially dispersed unmasking in the
Halton scheduler~\citep{besnier2025halton}. It also acts as a hard proxy for
the spatial-dispersion term in the choose-then-sample decomposition of
Proposition~4 in \citet{hayakawa2026demystifying}: confidence ordering favors
low-uncertainty coordinates, while spacing prevents one batch from being spent
entirely on nearby, potentially redundant positions. This connection motivates
the fixed rule rather than identifying an optimal radius for language.
The primary implementation distributes the path-total deficit budget and the
sequence-length count budget uniformly over the $R$ post-warmup steps:
\[
 E_r=D_{\max}/R,
 \qquad
 K_r=\lceil L/R\rceil,
 \qquad r=W,\ldots,N-2.
\]
The deficit budgets are upper bounds rather than quotas. In the primary
64-NFE duration factorial, $K_r=25$ after warmup and
$R_{\mathrm{sp}}=4$.
Smaller deficit means that the denoiser output places more mass on its preferred
label. The selector is a concentration controller rather than a calibrated
estimate of target-token error.

\begin{algorithm}[H]
\caption{Committed reveal with horizon $H$}
\label{alg:crs-horizon}
\small
\begin{algorithmic}[1]
\Require checkpoint $\mathsf M_\theta$; total NFE $N$ and grid
$t_0>\cdots>t_{N-1}=\varepsilon>0$;
warmup $W$; horizon $H\in\{0,1,\ldots,\infty\}$; selector budgets $E_r$ and
count caps $K_r$;
top-$p$ schedule $(p_r)_r$
\State sample $\bm x_{t_0}$ from the terminal uniform law
\State $C\gets\varnothing$; stored labels $\bm c\gets\varnothing$;
expiry map $e\gets\varnothing$; previously selected set $B\gets\varnothing$
\For{$r=0,\ldots,N-2$}
  \If{$0<H<\infty$}
    \State remove $i$ from $C$ whenever $r=e_i$
  \EndIf
  \State $\widetilde{\bm x}_r\gets\bm x_{t_r}[C\leftarrow\bm c|_C]$
  \State $h_r\gets\mathsf M_\theta(\widetilde{\bm x}_r,t_r)$
  \If{$r\ge W$ and $H>0$}
    \State choose $S_r$ from $[L]\setminus B$ using
    Equations~\eqref{eq:selector-deficit-app}--\eqref{eq:selector-constraints-app}
    \State $c_i\gets a_r^i$ for $i\in S_r$
    \If{$0<H<\infty$}
      \State $e_i\gets r+H+1$ for $i\in S_r$
    \EndIf
    \State $C\gets C\cup S_r$; $B\gets B\cup S_r$
  \EndIf
  \State independently sample each updated $x_{t_{r+1}}^i$ from
  Equation~\eqref{eq:overwrite-transition-app}
  \Comment{the transition uses $x_{t_r}^i$, not $\widetilde x_r^i$}
\EndFor
\If{$0<H<\infty$}
  \State remove $i$ from $C$ whenever $N-1=e_i$
\EndIf
\State use the final step at $t_{N-1}$ for the common argmax readout,
with input $\bm x_{t_{N-1}}[C\leftarrow\bm c|_C]$
\end{algorithmic}
\end{algorithm}

Newly stored labels do not alter the already computed transition at step $r$;
they first enter the denoiser at step $r+1$. With $e_i=r+H+1$, a label selected
at $r$ is visible for exactly the next $H$ steps (or all remaining steps if
fewer remain). For $H=0$, selection and storage are disabled and the algorithm
is the no-storage revisable path under the supplied top-$p$ schedule. For
$H=\infty$, no expiry occurs and the
algorithm is \CRS. In the finite-horizon factorial, an expired
coordinate becomes revisable but remains in $B$, so it cannot be selected
again. Thus $H=8$ is a finite-persistence control lasting eight later steps, not a
renewable policy.

\FloatBarrier

\section{Warmup as Lower-Risk Value Selection}
\label{app:warmup-theory}

\subsection{Proof of Proposition~\ref{prop:warmup-blackwell}}
\label{app:warmup-proposition-proof}

For $s<t$, coordinatewise factorization of the forward kernel gives the
Markov chain
\[
 X_0^i\longrightarrow\bm X_s^{-i}\longrightarrow\bm X_t^{-i}.
\]
Any decision rule based on $\bm X_t^{-i}$ can therefore be simulated from
$\bm X_s^{-i}$ by first drawing a synthetic later state from
$q_{t\mid s}^{-i}:=\prod_{j\ne i}q_{t\mid s}^j$, the product transition on
the non-$i$ coordinates. The joint law of that synthetic state and $X_0^i$ is the
same as the joint law of $\bm X_t^{-i}$ and $X_0^i$. Hence the best rule based
on $\bm X_s^{-i}$ cannot have larger zero--one risk. The Bayes rule is the
largest-probability LOO label, so
$\mathcal R_s^i\le\mathcal R_t^i$.

\subsection{Illustration: warmup separates sequence modes}
\label{app:warmup-codebook}

The proposition orders exact-law selection risks without assuming a particular
data distribution. To make its sequence-level consequence concrete, consider
$K$ candidate clean sequences $c_1,\ldots,c_K\in\mathcal S^L$ with minimum
Hamming separation
\[
 \Delta:=\min_{z\ne z'}d_H(c_z,c_{z'})>0.
\]
Let the mode $M$ be uniform on $[K]$, set $\bm X_0=c_M$, and condition on
$M=m_\star$, writing $c_\star:=c_{m_\star}$. For a forward-corrupted state
$\bm x_s$, put $d=d_H(\bm x_s,c_\star)$. Under the uniform channel,
\[
 q_{s\mid0}(\bm x_s\mid c_z)
 =v_s^{L-d_H(\bm x_s,c_z)}u_s^{d_H(\bm x_s,c_z)},
 \qquad
 v_s=1-\alpha_s+\frac{\alpha_s}{V},
 \quad u_s=\frac{\alpha_s}{V}.
\]
Because $v_s>u_s$ away from complete noise, equal-prior
maximum-a-posteriori (MAP) decoding is nearest-Hamming decoding. For every
competing $c_z$, the triangle inequality gives
\[
 d_H(\bm x_s,c_z)
 \ge d_H(c_z,c_\star)-d_H(\bm x_s,c_\star)
 \ge\Delta-d.
\]
Thus $d<\Delta/2$ makes $c_\star$ the unique nearest codeword and hence the
unique MAP mode.

This mode decision also connects to the coordinatewise argmax used by
\CRS{}. Let $w_z=\mathbb P(M=z\mid\bm X_s=\bm x_s)$, suppose
$w_\star:=w_{m_\star}>1/2$, and define
$R_w(a)=\sum_z w_zd_H(a,c_z)$. The reverse triangle inequality yields
\begin{equation}
 R_w(a)-R_w(c_\star)
 \ge(2w_\star-1)d_H(a,c_\star).
 \label{eq:hamming-bayes-gap-app}
\end{equation}
The right-hand side is positive for every $a\ne c_\star$, proving uniqueness
of the posterior-Hamming Bayes action. Because Hamming risk separates by
coordinate, this action is also obtained by taking the clean-token posterior
argmax at each coordinate.

The same model gives an explicit probability of entering this uniquely
decodable region. Let
\[
 r_s=\frac{(V-1)\alpha_s}{V}
\]
be the probability that one observed symbol differs from its clean symbol.
Then $d_H(\bm X_s,c_\star)\sim\operatorname{Binomial}(L,r_s)$. If
$r_s<\Delta/(2L)$, the standard binomial Chernoff bound gives
\begin{equation}
 \mathbb P\!\left(
 c_\star\text{ is not the unique nearest-Hamming mode}
 \mid c_\star
 \right)
 \le
 \exp\!\left\{
 -L\,d_{\mathrm{bin}}\!\left(
 \frac{\Delta}{2L}\,\middle\Vert\,r_s
 \right)
 \right\},
 \label{eq:warmup-codebook-bound-app}
\end{equation}
where
$d_{\mathrm{bin}}(a\,\Vert\,b)
=a\log(a/b)+(1-a)\log((1-a)/(1-b))$.
Thus lower noise both improves the general oracle-risk ordering in
Proposition~\ref{prop:warmup-blackwell} and, in this separated-mode example,
makes one coherent sequence mode more likely to be uniquely identifiable.
The codebook calculation supplies a concrete sufficient condition, while the
proposition itself provides the distribution-free exact-law comparison.

\FloatBarrier

\section{Shared-Mode Proof}
\label{app:shared-mode-proof}

\subsection{Proof of Theorem~\ref{thm:shared-mode}}
\label{app:shared-mode-theorem-proof}

We first make the conditional extension precise. Let $M\in[K]$ have
$\mathbb P(M=m)=w_m$ and suppose
\[
 \mathbb P(Y=y\mid M=m)
 =\prod_{i=1}^nP_{m,i}(y_i).
\]
Let $U$ be a potential visible-context observation and assume
$Y\perp\!\!\!\perp U\mid M$. For a realization $u$, define the context
likelihood $L_m(u):=\mathbb P(U=u\mid M=m)$. Whenever the denominator is
positive, Bayes' rule gives
\begin{equation}
 w_m[u]:=\mathbb P(M=m\mid U=u)
 =\frac{w_mL_m(u)}{\sum_{\ell=1}^K w_\ell L_\ell(u)}.
 \label{eq:context-posterior-weights-app}
\end{equation}
The conditional-independence assumption is essential: without it,
conditioning on arbitrary context can create dependence among the $Y_i$ even
when the mode is fixed.

The one-coordinate marginals of $Q_{\mathrm{mix}}^w$ are
$Q_{\mathrm{mix},i}^w(y_i)=\sum_m w_mP_{m,i}(y_i)$, hence
$Q_{\mathrm{prod}}^w=\prod_iQ_{\mathrm{mix},i}^w$. Therefore
\begin{equation}
 \KL(Q_{\mathrm{mix}}^w\kldiv Q_{\mathrm{prod}}^w)
 =\sum_i\Ent_{Q_{\mathrm{mix}}^w}(Y_i)
  -\Ent_{Q_{\mathrm{mix}}^w}(Y)
 =:\TC_{Q_{\mathrm{mix}}^w}(Y_1,\ldots,Y_n).
\end{equation}
Using the chain rule and conditional independence given $M$,
\begin{align}
 \TC_{Q_{\mathrm{mix}}^w}(Y_1,\ldots,Y_n)
 &=\sum_{i=2}^n I(Y_i;Y_{<i})
 \le\sum_{i=2}^n I(Y_i;M)
 \le(n-1)\Ent(M)=(n-1)\Ent(w).
\end{align}
Equation~\eqref{eq:context-posterior-weights-app} and
$Y\perp\!\!\!\perp U\mid M$ show that, after observing $U=u$, the law
of $Y$ has the same mixture-of-products form with $w$ replaced by
$w[u]$. Thus $Q_{\mathrm{mix}}^{w[u]}$ and
$Q_{\mathrm{prod}}^{w[u]}$ are precisely
the two laws in Equation~\eqref{eq:shared-mode-laws} evaluated at these
posterior weights. Repeating the argument gives the conditional total-correlation
bound. If $\epsilon(u)=1-\max_m w_m[u]$, separating the largest mode from
the remaining $K-1$ modes yields
\[
 \Ent(w[u])
 \le h_2(\epsilon(u))+\epsilon(u)\log(K-1),
\]
which proves Equation~\eqref{eq:context-mode-gap}.

For completeness, a simple likelihood-ratio condition makes the concentration
explicit. Fix a reference mode $m^\star$ with $w_{m^\star}>0$ and
$L_{m^\star}(u)>0$, and put
$\rho(u):=\max_{m\ne m^\star}L_m(u)/L_{m^\star}(u)$. Summing posterior odds in
Equation~\eqref{eq:context-posterior-weights-app} gives
\begin{equation}
 \frac{1-w_{m^\star}[u]}{w_{m^\star}[u]}
 \le\frac{1-w_{m^\star}}{w_{m^\star}}\rho(u).
 \label{eq:context-posterior-odds-app}
\end{equation}
Independent mode-consistent context observations multiply their likelihood
ratios, so a uniform ratio below one makes this bound decay exponentially in
the number of observations.
This result quantifies coordination around the mode favored by the observed
context, not whether that mode is the correct one. It also treats $U$ as an
exogenous observation satisfying the stated conditional-independence model;
the closed loop created when a sampler inserts its own argmax is assessed by
the empirical controls in Appendix~\ref{app:mechanism-tables}.

\subsection{Support evidence and persistent evidence}
\label{app:combined-evidence}

The theorem treats visible labels as evidence about a shared mode. We now place
one-step support restriction and later persistent context in the same posterior
calculation. Condition throughout on the current reverse history, so the
retained sets below are fixed. We use $U_{1:H}$ as an additive-evidence
abstraction for the step-specific states encountered while selected context
remains visible; it is not a model of independently observing the same stored
token $H$ times. Let $M\in[K]$ have prior weights $w_m$. Conditional on $M=m$,
suppose a support-bearing vector $X=(X_1,\ldots,X_d)$ has product law
$\prod_iR_{m,i}$, a future vector $Y=(Y_1,\ldots,Y_n)$ has product law
$\prod_jP_{m,j}$, and visible-context observations $U_1,\ldots,U_H$ are
mutually independent. Assume that $X$, $Y$, and $(U_e)_{e=1}^H$ are mutually
independent given $M$.

Fix a reference mode $m^\star$ and a retained rectangle $A=\prod_iA_i$. For
$m\ne m^\star$ and a realization $u_{1:H}$, define
\begin{align}
 S_m(A)
 &:=\sum_i\log\frac{R_{m^\star,i}(A_i)}{R_{m,i}(A_i)},
 &
 T_m(u_{1:H})
 &:=\sum_{e=1}^H
 \log\frac{\mathbb P(U_e=u_e\mid M=m^\star)}
 {\mathbb P(U_e=u_e\mid M=m)}.
 \label{eq:combined-evidence-parts-app}
\end{align}
Assume that these likelihood ratios are finite and positive, and put
\begin{equation}
 S(A):=\min_{m\ne m^\star}S_m(A),
 \qquad
 T(u_{1:H}):=\min_{m\ne m^\star}T_m(u_{1:H}).
 \label{eq:combined-evidence-minima-app}
\end{equation}

\begin{theorem}[Combined mode evidence]
\label{thm:combined-evidence}
Let $w[A,u]$ be the posterior mode law after observing $X\in A$ and
$U_{1:H}=u_{1:H}$. Then
\begin{equation}
 \frac{1-w_{m^\star}[A,u]}{w_{m^\star}[A,u]}
 \le
 \frac{1-w_{m^\star}}{w_{m^\star}}
 e^{-(S(A)+T(u_{1:H}))}.
 \label{eq:combined-odds-app}
\end{equation}
With
\begin{equation}
 a:=\frac{1-w_{m^\star}}{w_{m^\star}},
 \qquad
 \varepsilon(S,T):=\frac{ae^{-(S+T)}}{1+ae^{-(S+T)}},
 \qquad
 \bar\varepsilon:=\min\{\varepsilon(S(A),T(u_{1:H})),1-1/K\},
 \label{eq:combined-error-app}
\end{equation}
we have
$1-w_{m^\star}[A,u]\le\varepsilon(S(A),T(u_{1:H}))$ and
\begin{equation}
 \Ent(w[A,u])
 \le h_2(\bar\varepsilon)+\bar\varepsilon\log(K-1).
 \label{eq:combined-entropy-app}
\end{equation}
If $Q_{\mathrm{mix}}^{w[A,u]}$ is the mixture of the product laws of $Y$ under
$w[A,u]$ and $Q_{\mathrm{prod}}^{w[A,u]}$ is the product of its coordinate
marginals, then
\begin{equation}
 \KL(Q_{\mathrm{mix}}^{w[A,u]}\kldiv Q_{\mathrm{prod}}^{w[A,u]})
 \le(n-1)
 \left[h_2(\bar\varepsilon)+\bar\varepsilon\log(K-1)\right].
 \label{eq:combined-gap-app}
\end{equation}
\end{theorem}

\begin{proof}
For every $m\ne m^\star$, Bayes' rule and conditional independence give
\[
 \frac{w_m[A,u]}{w_{m^\star}[A,u]}
 =\frac{w_m}{w_{m^\star}}
 \exp\{-S_m(A)-T_m(u_{1:H})\}
 \le\frac{w_m}{w_{m^\star}}
 e^{-(S(A)+T(u_{1:H}))}.
\]
Summing proves Equation~\eqref{eq:combined-odds-app}, and solving the odds
inequality gives the posterior-error bound. When that upper bound is below
$1-1/K$, entropy is maximized by placing the remaining mass uniformly over
the other $K-1$ modes. For a larger upper bound, the uniform mode law is
feasible and the entropy certificate saturates at $\log K$. This proves
Equation~\eqref{eq:combined-entropy-app}. Applying
Theorem~\ref{thm:shared-mode} to the future product components gives
Equation~\eqref{eq:combined-gap-app}.
\end{proof}

Equation~\eqref{eq:combined-odds-app} also gives the mode-log-loss
certificate
\begin{equation}
 -\log w_{m^\star}[A,u]
 \le
 \log\!\left(1+a e^{-\{S(A)+T(u_{1:H})\}}\right)
 =:\mathcal L_{\mathrm{mode}}\!\left(S(A),T(u_{1:H})\right).
 \label{eq:mode-logloss-certificate-app}
\end{equation}

\begin{corollary}[Diminishing returns under additive mode evidence]
\label{cor:mode-diminishing-returns}
Assume $0<w_{m^\star}<1$. Write
$a=(1-w_{m^\star})/w_{m^\star}$ and use
$\mathcal L_{\mathrm{mode}}$ from
Equation~\eqref{eq:mode-logloss-certificate-app}.
For $\delta>0$, let the gain in this mode-log-loss certificate be
$\mathcal G_\delta(S,T):=\mathcal L_{\mathrm{mode}}(S,T)
-\mathcal L_{\mathrm{mode}}(S,T+\delta)$. Then
\begin{equation}
 \mathcal G_\delta(S,T)>0,
 \qquad
 \partial_S\mathcal G_\delta(S,T)
 =\partial_T\mathcal G_\delta(S,T)<0,
 \label{eq:mode-diminishing-returns}
\end{equation}
and $\mathcal G_\delta(S,T)\to0$ as $S+T\to\infty$.
\end{corollary}

\begin{proof}
Put $u=S+T$. Direct substitution gives
\[
 \mathcal G_\delta(S,T)
 =\log\frac{1+ae^{-u}}{1+ae^{-(u+\delta)}}>0.
\]
Because the gain depends on $S$ and $T$ only through their sum, its two partial
derivatives are equal. They are
\begin{equation}
 \partial_S\mathcal G_\delta(S,T)
 =\partial_T\mathcal G_\delta(S,T)
 =-\frac{ae^{-u}}{1+ae^{-u}}
 +\frac{ae^{-(u+\delta)}}{1+ae^{-(u+\delta)}}<0.
\end{equation}
The inequality follows because $x/(1+x)$ is strictly increasing. The displayed
ratio defining $\mathcal G_\delta$ tends to one as $u\to\infty$, so the gain
tends to zero.
\end{proof}

Under uniform corruption, if $g_{m,j}$ is the clean-token law at visible
coordinate $j$, token $c_j$ supplies Bayes factor
\begin{equation}
 B_{t,j}^{m^\star,m}(c_j)
 =\frac{(1-\alpha_t)g_{m^\star,j}(c_j)+\alpha_t/V}
 {(1-\alpha_t)g_{m,j}(c_j)+\alpha_t/V}.
 \label{eq:uniform-anchor-bf-app}
\end{equation}
When $g_{m^\star,j}(c_j)>g_{m,j}(c_j)$, this factor increases as
$\alpha_t$ decreases. A mode-consistent visible token therefore becomes more
informative toward the clean end of the reverse path.

\iffalse
\section{Adaptive Current-Step Operators}
\label{app:adaptive-current-step}

The denoiser-output shape scores used in the adaptive audit admit a simple ordering.
Let $q_1\ge q_2$ be the top two clean probabilities, put
$s=q_1+q_2$ and $m=(q_1-q_2)/(q_1+q_2)$, and define
\begin{align}
 c_{\mathrm{norm}}&:=m,
 &c_{\mathrm{abs}}&:=sm=q_1-q_2,\nonumber\\
 c_{\mathrm{bin}}&:=s\left[1-
 \frac{h_2((1+m)/2)}{\log2}\right].
 \label{eq:score-definitions-app}
\end{align}

\begin{lemma}[Ordering of top-two certainty scores]
\label{lem:score-ordering}
For $q_1\ge q_2\ge0$ with $q_1+q_2>0$,
\begin{equation}
 0\le c_{\mathrm{bin}}
 \le c_{\mathrm{abs}}
 \le c_{\mathrm{norm}}
 \le1.
 \label{eq:score-ordering-app}
\end{equation}
\end{lemma}

\begin{proof}
Concavity of $h_2$ on $[1/2,1]$ places it above the chord joining
$(1/2,\log2)$ and $(1,0)$. Thus
$h_2((1+m)/2)\ge(1-m)\log2$, so the bracket in
$c_{\mathrm{bin}}$ lies in $[0,m]$. Multiplying by $s\in(0,1]$ gives
Equation~\eqref{eq:score-ordering-app}.
\end{proof}

\begin{theorem}[A charge-matched interface certificate]
\label{thm:charge-matched-certificate}
Assume $0<w_{m^\star}<1$. Consider two restriction allocations $W$ and $G$ at
one reverse state, with
$\sum_i\chi_i^W=\sum_i\chi_i^G$. Suppose their retained rectangles have
shared-mode evidence $S_W>S_G$ and their charge vectors satisfy
$\boldsymbol\chi^W\succ\boldsymbol\chi^G$. Then, for every fixed amount of
subsequent compatible context evidence $T$,
\begin{equation}
 \mathcal L_{\mathrm{mode}}(S_W,T)
 <\mathcal L_{\mathrm{mode}}(S_G,T),
 \qquad
 \sum_i\Phi_V(\chi_i^W)
 \le\sum_i\Phi_V(\chi_i^G).
 \label{eq:charge-matched-certificate-app}
\end{equation}
\end{theorem}

\begin{proof}
The first comparison follows because
$\mathcal L_{\mathrm{mode}}(S,T)=\log(1+ae^{-(S+T)})$ is strictly decreasing
in $S$. The second is Corollary~\ref{cor:charge-majorization}.
\end{proof}

This is an interface-level sufficient condition. It concerns a mode-log-loss
certificate and a worst-case sum of marginal entropy disturbances, rather than
an endpoint evaluator or diversity frontier.

Charge and marginal shape still do not identify which coordinate is relevant
to a shared mode.

\begin{proposition}[Marginal charge does not determine mode relevance]
\label{prop:charge-no-free-lunch}
When a budget permits restricting one of two coordinates, no rule based only on
their coordinate marginals and restriction charges can always select the more
informative coordinate for identifying a shared latent mode.
\end{proposition}

\begin{proof}
Let $Z\sim\operatorname{Bernoulli}(1/2)$. In model $\mathcal M_1$, let $X_1$
equal $Z$ with probability $1-\epsilon$ and be flipped otherwise, while
$X_2\sim\operatorname{Bernoulli}(1/2)$ is independent of $Z$. In model
$\mathcal M_2$, exchange the roles of $X_1$ and $X_2$. Both coordinates have
the same $\operatorname{Bernoulli}(1/2)$ marginal in both models, and
restricting either to $\{0\}$ has charge $\log2$. In $\mathcal M_1$, the first
event supplies mode log likelihood ratio
$\log((1-\epsilon)/\epsilon)$ while the second supplies zero; the ordering
reverses in $\mathcal M_2$.
\end{proof}

Thus $\chi$ is an exact intervention ledger, while the score $s_i$ is an
explicit modeling choice about relevance. The allocation controls in
Appendix~\ref{app:mechanism-tables} test that choice by preserving, permuting,
flattening, or reversing the score assignment.
\fi

\FloatBarrier

\section{Experimental Details}
\label{app:experiment-details}

\subsection{Pathwise singleton trace}
\label{app:singleton-trace}

We trace the native Duo-distilled sampler at 64 NFE using four seed shards and
four paths per shard. To make the trace quantities explicit, at reverse step
$r$ define
\begin{equation}
 h_r^i(a):=h_{\theta,t_r}^i(a\mid\bm x_{t_r}),
 \qquad
 A_{r,i}^{(p)}:=A_p(h_r^i),
 \qquad
 a_r^i:=\arg\max_a h_r^i(a).
 \label{eq:singleton-trace-def-app}
\end{equation}
A singleton event is
$E_{r,i}^{(p)}:=\{A_{r,i}^{(p)}=\{a_r^i\}\}$. On this event,
$\mathsf T_ph_r^i=\delta_{a_r^i}$ and Bayes conversion preserves that point
mass. Hence the transition sampled by the current stochastic bridge reduces to
\begin{equation}
 X_{t_{r+1}}^i\sim
 q_{t_{r+1}\mid0,t_r}^i
 \bigl(\,\cdot\mid a_r^i,x_{t_r}^i\bigr),
 \label{eq:singleton-bridge-app}
\end{equation}
which need not be a point mass even though the clean-token proposal is one.

The table reports averages of the following event indicators. ``Next step''
is
\begin{equation}
 \1\!\left\{
 \arg\max_a h_{r+1}^i(a)=a_r^i
 \right\}
 \quad\text{conditioned on }E_{r,i}^{(p)}
 \text{ and on step }r+1\text{ existing},
 \label{eq:singleton-next-step-app}
\end{equation}
so it asks whether the point hypothesis remains the denoiser's preferred clean
label after one sampled bridge transition; it does not require the next top-$p$
support to remain a singleton. ``Final'' is
\begin{equation}
 \1\{\widehat x^i=a_r^i\},
 \qquad
 \widehat x^i=\arg\max_a
 h_{\theta,\varepsilon}^i(a\mid\bm x_\varepsilon),
 \label{eq:singleton-final-app}
\end{equation}
conditioned on $E_{r,i}^{(p)}$. ``Stochastic bridge'' is the conditional rate
of
\begin{equation}
 \max_y q_{t_{r+1}\mid0,t_r}^i
 (y\mid a_r^i,x_{t_r}^i)<1-10^{-8},
 \label{eq:singleton-stochastic-bridge-app}
\end{equation}
which distinguishes a point clean hypothesis from a deterministic update of
the mutable noisy state. ``Singleton'' averages $\1\{E_{r,i}^{(p)}\}$ over all
active coordinate--step pairs; ``Last 10'' restricts that average to the final
ten reverse steps.

% Numeric display policy for empirical tables: GenPPL/PPL uses two decimals,
% unigram entropy uses three decimals, and bounded diversity proportions use
% four decimals. Paired contrasts retain the extra precision needed to resolve
% their reported SEs; externally reported values retain their source precision.
\begin{table}[H]
\centering
\small
\setlength{\tabcolsep}{2.8pt}
\begin{tabular}{@{}lccccc@{}}
\toprule
Native $p$ & Singleton & Last 10 & Next step & Final & Stochastic bridge \\
\midrule
0.8 & $35.44\!\pm\!0.32$ & $58.46\!\pm\!0.90$ & $99.65\!\pm\!0.03$ &
$95.93\!\pm\!0.20$ & $94.80\!\pm\!0.06$ \\
0.9 & $27.54\!\pm\!0.43$ & $45.56\!\pm\!0.24$ & $99.78\!\pm\!0.02$ &
$97.16\!\pm\!0.12$ & $94.76\!\pm\!0.09$ \\
\bottomrule
\end{tabular}
\caption{Pathwise point hypotheses, in percent (mean$\pm$standard error (SE)
over four seed
shards). `Last 10' is the singleton rate over the final ten reverse
transitions.}
\label{tab:singleton-trace-app}
\end{table}

\subsection{NFE scaling on the two Duo checkpoints}
\label{app:nfe-scaling}

The common-decoder NFE study uses four aligned seed shards and 64 samples per
cell per shard. Both public Duo checkpoints generate 1024 GPT-2 tokens with
path-total $D_{\max}=512$ for CRS. In both cases the warmup fraction is $0.33$,
CRS uses full clean-token support and persistent input visibility ($H=\infty$), and every
procedure uses argmax decoding in the final step. Initial states and
random-number-generator (RNG) streams
are reset within each NFE. The NFE sweep sets
$K_r=\lceil1024/(N-1-W)\rceil$, giving per-step caps $205,103,49,25$ at
$N\in\{8,16,32,64\}$ and approximately 1024 positions of path capacity after
warmup. The subtraction by one reserves the final step for the
common argmax readout rather than a reverse transition.

At 64 NFE, the first selection follows 21 warmup steps. With
$D_{\max}=512$, CRS cumulatively selects $991.23\pm1.02$ coordinates per sample
on Duo-distilled and $963.73\pm0.73$ on Duo. Across the 42 post-warmup steps,
the mean number of simultaneously visible stored labels is
$508.23\pm0.19$ and $502.68\pm0.16$, respectively. Values are mean$\pm$SE over
the four seed shards. These path statistics show that the intervention is
gradually accumulated rather than a one-shot overwrite of the sequence.

\begin{table}[H]
\centering
\small
\setlength{\tabcolsep}{4.2pt}
\begin{tabular}{@{}lrrrr@{}}
\toprule
Checkpoint / NFE & Full support & top-$p$ 0.95 & top-$p$ 0.9 & CRS \\
\midrule
Duo-distilled / 8  & 76.89 / 5.269 & 65.60 / 5.282 & 56.11 / 5.255 & 39.57 / 5.136 \\
Duo-distilled / 16 & 54.58 / 5.376 & 44.91 / 5.362 & 38.82 / 5.324 & 31.80 / 5.318 \\
Duo-distilled / 32 & 46.76 / 5.380 & 37.94 / 5.360 & 32.57 / 5.320 & 30.08 / 5.338 \\
Duo-distilled / 64 & 41.79 / 5.351 & 34.34 / 5.326 & 29.76 / 5.296 & 28.27 / 5.317 \\
\midrule
Duo / 8  & 131.66 / 4.897 & 109.31 / 4.938 & 95.32 / 4.943 & 46.04 / 4.755 \\
Duo / 16 &  99.31 / 5.138 &  78.18 / 5.189 & 64.41 / 5.176 & 37.49 / 5.105 \\
Duo / 32 &  85.11 / 5.216 &  65.72 / 5.242 & 54.66 / 5.229 & 37.26 / 5.185 \\
Duo / 64 &  78.31 / 5.231 &  59.53 / 5.250 & 47.20 / 5.220 & 35.75 / 5.184 \\
\bottomrule
\end{tabular}
\caption{NFE scaling under common argmax decoding in the final step. Each
entry is GenPPL / mean within-sequence unigram entropy over 256 generations.
The architecture, tokenizer, and sequence length are shared; the checkpoint is
the only model-level change across the horizontal rule. The dense base-Duo
64-NFE frontier in Table~\ref{tab:frontier-values-app} is an independent launch
rather than a pooled estimate.}
\label{tab:nfe-scaling-app}
\end{table}

\begin{table}[H]
\centering
\small
\setlength{\tabcolsep}{3.5pt}
\begin{tabular}{@{}lrrrr@{}}
\toprule
Checkpoint / NFE & Full support & top-$p$ 0.95 & top-$p$ 0.9 & CRS \\
\midrule
Duo-distilled / 8  & 0.84 / 0.006 & 0.79 / 0.005 & 0.31 / 0.008 & 0.11 / 0.005 \\
Duo-distilled / 16 & 0.43 / 0.001 & 0.26 / 0.004 & 0.18 / 0.006 & 0.26 / 0.004 \\
Duo-distilled / 32 & 0.64 / 0.007 & 0.59 / 0.010 & 0.51 / 0.009 & 0.18 / 0.008 \\
Duo-distilled / 64 & 0.84 / 0.015 & 0.39 / 0.014 & 0.19 / 0.009 & 0.53 / 0.015 \\
\midrule
Duo / 8  & 1.61 / 0.003 & 0.99 / 0.003 & 1.15 / 0.002 & 0.19 / 0.012 \\
Duo / 16 & 0.76 / 0.009 & 0.56 / 0.003 & 0.64 / 0.006 & 0.12 / 0.004 \\
Duo / 32 & 1.62 / 0.008 & 1.40 / 0.011 & 0.90 / 0.010 & 0.64 / 0.012 \\
Duo / 64 & 1.62 / 0.011 & 0.64 / 0.009 & 0.77 / 0.009 & 0.62 / 0.012 \\
\bottomrule
\end{tabular}
\caption{Four-shard standard errors for Table~\ref{tab:nfe-scaling-app}.
Each entry is GenPPL SE / unigram-entropy SE. Figure~\ref{fig:nfe-scaling}
shows point estimates without bars for legibility.}
\label{tab:nfe-scaling-se-app}
\end{table}

\paragraph{Paired endpoint audit.}
\mbox{\citet{deschenaux2026psi}} report the final
noise-removal step as an ancestral categorical draw, whereas our main comparison
uses coordinatewise argmax decoding in the final step. To isolate this protocol
difference, Table~\ref{tab:deschenaux-endpoint-audit} reuses each pre-terminal
reverse path and changes only that final output rule. The ancestral base-Duo cells recover the published 32-NFE
operating region at both tested support levels. The effect is not a uniform
quality bonus: argmax lowers entropy in every row, improves native/full-support
GenPPL, and worsens base-Duo GenPPL at $p=0.9$. Endpoint choice therefore
substantially accounts for the apparent numerical discrepancy, but its effect
depends on checkpoint and support restriction.

\begin{table}[H]
\centering
\small
\setlength{\tabcolsep}{3.8pt}
\begin{tabular}{@{}llrrr@{}}
\toprule
Checkpoint & Support & Reported ancestral & Paired ancestral & Paired argmax \\
\midrule
Duo & $p=1.0$ & $96.76\;(5.57)$ & $100.11\;(5.58)$ & $87.62\;(5.23)$ \\
Duo & $p=0.9$ & $44.24\;(5.40)$ & $44.78\;(5.40)$ & $53.57\;(5.22)$ \\
Duo-distilled & $p=1.0$ & $68.35\;(5.54)$ & $61.13\;(5.54)$ & $45.50\;(5.40)$ \\
Duo-distilled & $p=0.9$ & $35.92\;(5.41)$ & $34.90\;(5.42)$ & $33.54\;(5.35)$ \\
\bottomrule
\end{tabular}
\caption{32-NFE endpoint audit. Entries are GenPPL (within-sequence unigram
entropy, shown to the two-decimal precision available in the reported source).
The two paired columns share the pre-terminal path and differ only in
the last output rule; each uses four seed shards and 64 generations. Values in
the reported column are read from the ancestral-sampling tables of
\mbox{\citet{deschenaux2026psi}}. These paired 64-generation cells are
independent of the 256-generation NFE sweep in
Table~\ref{tab:nfe-scaling-app}, so their argmax point estimates need not
coincide.}
\label{tab:deschenaux-endpoint-audit}
\end{table}

\begin{table}[H]
\centering
\footnotesize
\begin{minipage}[t]{0.48\linewidth}
\centering
\textbf{Duo-distilled}\\[2pt]
\begin{tabular}{@{}lrr@{}}
\toprule
Control & GenPPL & Entropy \\
\midrule
$p=1.0$ & $41.79\pm0.84$ & $5.351\pm0.015$ \\
$p=0.95$ & $34.34\pm0.39$ & $5.326\pm0.014$ \\
$p=0.9$ & $29.76\pm0.19$ & $5.296\pm0.009$ \\
$p=0.85$ & $25.94\pm0.34$ & $5.251\pm0.013$ \\
$p=0.8$ & $23.34\pm0.39$ & $5.205\pm0.016$ \\
$p=0.75$ & $20.38\pm0.23$ & $5.148\pm0.017$ \\
\midrule
$D=128$ & $33.09\pm0.69$ & $5.336\pm0.014$ \\
$D=256$ & $30.54\pm0.62$ & $5.328\pm0.015$ \\
$D=512$ & $28.27\pm0.53$ & $5.317\pm0.015$ \\
$D=1024$ & $27.62\pm0.42$ & $5.312\pm0.014$ \\
\bottomrule
\end{tabular}
\end{minipage}
\hfill
\begin{minipage}[t]{0.48\linewidth}
\centering
\textbf{Duo}\\[2pt]
\begin{tabular}{@{}lrr@{}}
\toprule
Control & GenPPL & Entropy \\
\midrule
$p=1.0$ & $80.35\pm2.28$ & $5.233\pm0.011$ \\
$p=0.95$ & $59.53\pm0.64$ & $5.250\pm0.009$ \\
$p=0.9$ & $48.42\pm1.82$ & $5.221\pm0.014$ \\
$p=0.85$ & $39.86\pm0.73$ & $5.195\pm0.009$ \\
$p=0.8$ & $32.34\pm0.72$ & $5.114\pm0.009$ \\
$p=0.75$ & $28.33\pm0.64$ & $5.075\pm0.012$ \\
\midrule
$D=128$ & $51.96\pm0.78$ & $5.243\pm0.011$ \\
$D=256$ & $41.99\pm0.46$ & $5.216\pm0.011$ \\
$D=512$ & $35.75\pm0.62$ & $5.184\pm0.012$ \\
$D=1024$ & $34.60\pm0.53$ & $5.170\pm0.014$ \\
\bottomrule
\end{tabular}
\end{minipage}
\caption{Mean $\pm$ SE over four seed shards for the 64-NFE operating points.
All generations have length 1024. The base-Duo top-$p$ curve combines
protocol-aligned runs: $p=0.95$ comes from the fixed-procedure NFE sweep and
the remaining thresholds come from the denser 64-NFE sweep. Those denser reruns
give $80.35/48.42$ rather than $78.31/47.20$ GenPPL for $p=1.0/0.9$ in
Table~\ref{tab:nfe-scaling-app}; we do not pool the launches. The base-Duo
comparison is therefore descriptive.}
\label{tab:frontier-values-app}
\end{table}

The 8-NFE CRS cells commit $799.94$ positions on Duo-distilled and $795.50$ on
Duo on average. With a one-third warmup, both begin selection after two steps.
The frozen path-total budget is therefore aggressive relative to the available
low-NFE trajectory. These points diagnose one fixed procedure under reduced
compute rather than a tuned low-NFE frontier.

The base-Duo top-$p$ curve folds between $p=1.0$ and $p=0.95$ in unigram
entropy. We therefore do not invert that segment. Quantitative matched checks
use the monotone $p\le0.9$ branch and $D_{\max}\in\{256,512,1024\}$; the
$D_{\max}=128$ point is displayed only as a descriptive endpoint.

\subsection{How do the checkpoints differ in selection?}
\label{app:checkpoint-selector}

We evaluate both checkpoints on the same exact uniform-forward states and let
each checkpoint select 25 positions by Top-1 confidence with radius-four
exclusion. Both checkpoints are then evaluated on both selected sets. Across
$t\in\{0.7,0.67,0.6,0.5,0.4,0.2\}$, Duo-distilled self-selection has clean-label
accuracy $0.9$--$2.1$ percentage points higher than base-Duo self-selection. The selected-position Jaccard index
is only $0.43$--$0.53$. In contrast, when the selected positions are held fixed,
the two checkpoints agree on the argmax label at least $99.75\%$ of the time,
and their clean-label accuracy differs by at most $0.25$ points. Thus, in this
shared-state audit, the checkpoints differ mainly in which positions they
select. Once those positions are fixed, their argmax labels and clean-label
accuracies are nearly identical. Both selectors remain
overconfident: mean confidence is nearly one while selected-token error remains
$3.5$--$8.1\%$. This shared-state diagnostic does not include writeback or a
generated trajectory, so it does not attribute the full free-running endpoint
difference.

\subsection{Sampler-level warmup-by-duration ablation}
\label{app:warmup-duration}

We vary the full-support warmup fraction and the visibility horizon on
Duo-distilled at 64 NFE. The warmup fractions
$0.016,0.17,0.33,0.50$ correspond to $1,11,21,32$ warmup steps. Every
persistent cell uses $D_{\max}=512$. All cells retain full clean-token support
at every step and use argmax decoding in the final step. The cells share pair
seeds and RNG resets within each shard. Changing the requested warmup also
changes the two-piece reverse time grid, even when $H=0$ disables selection;
the $H=0$ rows therefore control for this schedule split rather than producing
bitwise-identical paths. Four paired seed shards contribute 256 generations
per cell.

\begin{table}[H]
\centering
\footnotesize
\setlength{\tabcolsep}{4.0pt}
\begin{tabular}{@{}crrrr@{}}
\toprule
Warmup steps & Native GenPPL & CRS GenPPL & Native entropy & CRS entropy \\
\midrule
1  & $42.15\pm0.98$ & $20.26\pm0.18$ & $5.359\pm0.014$ & $5.198\pm0.013$ \\
11 & $41.87\pm0.78$ & $25.93\pm0.27$ & $5.357\pm0.015$ & $5.297\pm0.013$ \\
21 & $41.79\pm0.84$ & $28.27\pm0.53$ & $5.351\pm0.015$ & $5.317\pm0.015$ \\
32 & $42.18\pm0.74$ & $29.83\pm0.47$ & $5.356\pm0.014$ & $5.328\pm0.014$ \\
\bottomrule
\end{tabular}
\caption{Endpoint means $\pm$ shard SE for the two policies plotted in
Figure~\ref{fig:warmup-permanent}. Native sampling keeps every position
revisable; \CRS{} keeps each selected label visible for all remaining
steps.}
\label{tab:warmup-permanent-endpoints-app}
\end{table}

\begin{table}[H]
\centering
\footnotesize
\setlength{\tabcolsep}{4.0pt}
\begin{tabular}{@{}crrrr@{}}
\toprule
Warmup steps & $\Delta$GenNLL $(8-0)$ & $\Delta$entropy $(8-0)$ &
$\Delta$GenNLL $(\infty-8)$ & $\Delta$entropy $(\infty-8)$ \\
\midrule
1  & $-0.486\pm0.023$ & $-0.1276\pm0.0082$ & $-0.246\pm0.004$ & $-0.0336\pm0.0022$ \\
11 & $-0.265\pm0.010$ & $-0.0486\pm0.0023$ & $-0.214\pm0.005$ & $-0.0123\pm0.0037$ \\
21 & $-0.211\pm0.008$ & $-0.0343\pm0.0020$ & $-0.180\pm0.003$ & $+0.0004\pm0.0013$ \\
32 & $-0.202\pm0.006$ & $-0.0335\pm0.0011$ & $-0.144\pm0.006$ & $+0.0052\pm0.0015$ \\
\bottomrule
\end{tabular}
\caption{Paired longer-horizon minus shorter-horizon contrasts (mean$\pm$SE
over four seed shards). Negative values mean lower evaluator GenNLL or lower
within-sequence unigram entropy under longer visibility.}
\label{tab:warmup-duration-app}
\end{table}

The $H=0$ endpoint is stable under the artificial phase split: GenPPL ranges
from $41.79$ to $42.18$, and unigram entropy ranges from $5.351$ to $5.359$.
This stability localizes the interaction to persistent exposure rather than
the time-grid split.
Delaying selection reduces both the likelihood gain and the entropy cost of
persistent context. At warmup fraction $0.33$, extending $H=8$ to $H=\infty$
leaves unigram entropy nearly unchanged, but the distinct 4-gram
fraction (Dist-4)~\citep{li2016diversity} decreases from $0.9405$ to $0.9359$, the repeated 4-gram
fraction (Rep-4) increases from
$0.0330$ to $0.0367$, and the across-sample similarity score
Self-BLEU~\citep{papineni2002bleu,zhu2018texygen}
increases from $0.2205$ to $0.2248$. The sequence-level movement
shows that unigram entropy understates the full duration cost. This
sampler-level ablation tests the qualitative warmup mechanism; quantitative
calibration of the oracle risk in Proposition~\ref{prop:warmup-blackwell} and
selection of an optimal warmup remain separate questions.

\subsection{Unconditional factorial}

We use `s-sahoo/duo-distilled' with the GPT-2 tokenizer. Every cell has 64 NFE,
length 1024, warmup fraction $0.33$, full-support warmup, canonical confidence
and spacing selection, $D_{\max}=512$, spacing radius four, float64 sampling
arithmetic to avoid the finite-precision categorical-sampling boundary effects
discussed by \citet{zheng2025secretly}, per-step count cap 25, and argmax decoding in the final step. Each of four seed shards
$8200,8210,8220,8230$ contributes 64 paired samples per cell, for 256 total. Initial states
and RNG streams are reset within each pair.
The released Duo top-$p$ implementation retains the first label beyond an
exact cumulative-mass equality; this floating-point boundary convention does
not affect the conversion identity in Section~\ref{sec:top-p-loo},
which conditions on the retained set actually used.

GenNLL is the mean next-token negative log likelihood under GPT-2-Large and
GenPPL is its exponential. Pythia-1.4B supplies the second evaluator. Entropy
is the mean, over generated sequences, of each sequence's empirical unigram
token entropy. It is distinct from an across-sequence distributional entropy.
For a multiset $\mathcal X$ of generated sequences and the multiset
$\mathcal G_n(x)$ of $n$-grams in $x$, our conventions are
\[
 \operatorname{Dist}\text{-}n
 =\frac{|\bigcup_{x\in\mathcal X}\operatorname{uniq}\mathcal G_n(x)|}
 {\sum_{x\in\mathcal X}|\mathcal G_n(x)|},
 \qquad
 \operatorname{Rep}\text{-}n
 =\frac1{|\mathcal X|}\sum_{x\in\mathcal X}
 \left(1-\frac{|\operatorname{uniq}\mathcal G_n(x)|}
 {|\mathcal G_n(x)|}\right).
\]
Self-BLEU-4 compares each
generation with the remaining generations. Higher entropy and Dist-4 indicate
broader lexical usage; lower Rep-4 and Self-BLEU indicate less repetition or
greater across-sample variation.
All entropy and NLL values use natural logarithms.
Because sentence-level BLEU depends on tokenization and
smoothing~\citep{post2018clarity,chen2014smoothing}, the unconditional table
uses leave-one-out GPT-2-tokenized Self-BLEU with equal weights over one- to
four-grams and add-one smoothing. The fixed-prefix analysis uses the same
weights with method-1 smoothing within each prompt group.

\begin{table}[H]
\centering
\small
\setlength{\tabcolsep}{4pt}
\begin{tabular}{@{}ccrrrrr@{}}
\toprule
Late $p$ & $H$ & GenPPL & Entropy & Dist-4 & Rep-4 & Self-BLEU \\
\midrule
1.0 & 0 & 41.79 & 5.351 & 0.9483 & 0.0272 & 0.2113 \\
1.0 & 8 & 33.84 & 5.317 & 0.9405 & 0.0330 & 0.2205 \\
1.0 & $\infty$ & 28.27 & 5.317 & 0.9359 & 0.0367 & 0.2248 \\
0.9 & 0 & 32.79 & 5.332 & 0.9432 & 0.0306 & 0.2207 \\
0.9 & 8 & 26.94 & 5.303 & 0.9346 & 0.0368 & 0.2294 \\
0.9 & $\infty$ & 26.40 & 5.302 & 0.9323 & 0.0389 & 0.2292 \\
0.85 & 0 & 30.58 & 5.321 & 0.9415 & 0.0316 & 0.2235 \\
0.85 & 8 & 25.66 & 5.293 & 0.9328 & 0.0388 & 0.2293 \\
0.85 & $\infty$ & 25.74 & 5.292 & 0.9328 & 0.0391 & 0.2285 \\
0.8 & 0 & 28.38 & 5.302 & 0.9388 & 0.0334 & 0.2275 \\
0.8 & 8 & 24.51 & 5.279 & 0.9308 & 0.0399 & 0.2331 \\
0.8 & $\infty$ & 24.89 & 5.283 & 0.9299 & 0.0414 & 0.2305 \\
\bottomrule
\end{tabular}
\caption{Complete unconditional top-$p$ by visibility-duration factorial. ``Late'' means
that the listed threshold is applied after full-support warmup. Means are over
four paired shards and 256 generations per cell.}
\label{tab:unconditional-full-app}
\end{table}

\subsection{Fixed-prefix transfer}

We use 1024 held-out OpenWebText~\citep{gokaslan2019openwebtext} prefixes, each 50 GPT-2 tokens long, with five
paired 50-token continuations per cell. The 50 prefix coordinates remain fixed
throughout sampling and are excluded from the selector; diversity summaries
use continuation tokens only. The sampler uses 64 NFE, the path-total
confidence-deficit budget $D_{\max}=25$, and argmax decoding in the final step. Reported conditional
perplexity (PPL) scores the continuation given the fixed prefix; conditional
negative log-likelihood (NLL)
is its log-scale counterpart. We evaluate 5120 generations per cell.

\begin{table}[H]
\centering
\small
\setlength{\tabcolsep}{3pt}
\begin{tabular}{@{}lcrrrrr@{}}
\toprule
Top-$p$ schedule & $H$ & GPT-2 & Pythia & Entropy & Dist-4 & Rep-4 \\
\midrule
Late 1.0 & 0 & 128.54 & 110.39 & 3.508 & 0.9606 & 0.0264 \\
Late 1.0 & 8 & 103.50 & 88.79 & 3.471 & 0.9480 & 0.0380 \\
Late 1.0 & $\infty$ & 83.96 & 71.75 & 3.453 & 0.9373 & 0.0469 \\
Late 0.9 & 8 & 83.21 & 71.62 & 3.461 & 0.9402 & 0.0446 \\
Late 0.8 & 0 & 88.43 & 76.71 & 3.488 & 0.9497 & 0.0356 \\
Late 0.8 & 8 & 74.77 & 64.55 & 3.452 & 0.9353 & 0.0482 \\
Late 0.8 & $\infty$ & 73.31 & 62.83 & 3.446 & 0.9341 & 0.0491 \\
\midrule
All 0.8 & 0 & 72.13 & 63.25 & 3.450 & 0.9393 & 0.0424 \\
\bottomrule
\end{tabular}
\caption{Fixed-prefix transfer point estimates. Entries labeled \emph{Late} use full-support
warmup and the listed $p$ afterward; \emph{All} uses that $p$ throughout.
GPT-2 and Pythia columns are conditional perplexity. Diversity summaries use
continuation tokens. Paired shard-level uncertainties for the primary
evaluator-NLL contrasts are reported below.}
\label{tab:conditional-full-app}
\end{table}

Under GPT-2-Large, the paired $H=0\to8$ NLL reductions are
$0.2165\pm0.0061$ at late $p=1.0$ and $0.1678\pm0.0053$ at late $p=0.8$;
the $H=8\to\infty$ reductions are $0.2091\pm0.0050$ and
$0.0198\pm0.0021$. The corresponding paired difference-in-differences are
$0.0487\pm0.0056$ and $0.1893\pm0.0069$ nats per token. Pythia-1.4B gives
interaction estimates $0.0451\pm0.0064$ and $0.1857\pm0.0093$. Positive
interaction values mean that lowering late $p$ attenuates the horizon effect.

For the paired evaluator-NLL contrasts against $p=0.8$ at every step with $H=0$, the
late-$p=0.8,H=8$ policy changes GPT-2/Pythia NLL by
$+0.0361\pm0.0105$/$+0.0208\pm0.0111$ nats per token; the $H=\infty$ policy
changes them by $+0.0163\pm0.0096$/$-0.0064\pm0.0100$. These are policy
contrasts on different generated trajectories and do not estimate the effect
of adding persistent context to the every-step schedule.
Within-prompt continuation Self-BLEU is $0.0378$ for every-step $p=0.8,H=0$,
$0.0351$ for late $p=0.8,H=8$, and $0.0357$ for late $p=0.8,H=\infty$.

\subsection{Support-by-duration interaction}
\label{app:support-duration}

Figure~\ref{fig:space-time} contains a complete unconditional factorial crossing
post-warmup top-$p\in\{1.0,0.9,0.85,0.8\}$ with visibility duration
$H\in\{0,8,\infty\}$. Temporary eight-step context remains useful throughout
that sweep. The marginal gain from extending $H=8$ to $H=\infty$
rapidly disappears as support restriction becomes stronger and becomes
slightly negative at late $p=0.8$. Fixed-prefix continuation has complete
duration contrasts only at $p\in\{1.0,0.8\}$. Its measured long-horizon gain
is also much smaller at $p=0.8$ than at $p=1.0$, although it remains positive;
these two endpoints do not constitute the full factorial.

\begin{figure}[H]
  \centering
  \includegraphics[width=\linewidth]{figures/space_time_interaction.pdf}
  \caption{Interaction between post-warmup support restriction and context
  duration. The unconditional series is the complete $4\times3$ factorial.
  The left panel shows the gain from keeping selected context visible
  for eight steps. The right panel shows the additional gain from keeping it
  visible for all remaining steps. Blue lines are unconditional generation; orange
  squares are the two observed fixed-prefix endpoints at $p=1.0$ and $p=0.8$.
  Moving right means stronger
  restriction. Error bars are standard errors over four paired shards.}
  \label{fig:space-time}
\end{figure}

Corollary~\ref{cor:mode-diminishing-returns} motivates this qualitative sign
test. We do not identify its idealized support evidence $S$ or context evidence
$T$ with a checkpoint statistic. The factorial instead asks whether stronger
support restriction attenuates the measured marginal horizon gain. The
corollary assumes exact, mode-consistent evidence, whereas the experiment uses
a learned denoiser on its own free-running trajectory. The slight negative
unconditional endpoint therefore lies outside the theorem's guaranteed regime
rather than contradicting its exact-model statement.

\subsection{Earlier storage-output control}
\label{app:storage-output-control}

This earlier 64-NFE control uses a $p=0.9$ warmup and sweeps post-warmup
top-$p$. Final-only and CRS store the same selected labels and copy them to the
same returned coordinates; they differ in whether those labels are also
inserted into later denoiser inputs. Their paired difference isolates repeated
input visibility beyond output storage. The fully revisable arm instead uses
argmax decoding in the final step and is included as a reference, rather than
as an output-rule-matched third arm.

\begin{figure}[H]
  \centering
  \includegraphics[width=0.9\linewidth]{figures/duo_distilled_common_endpoint_topp.pdf}
  \caption{Earlier storage--feedback control. Panel (a) uses unigram entropy
  as the diversity coordinate; panel (b) shows the post-warmup top-$p$ sweep.
  Vertical error bars show GenPPL standard errors over four paired shards; at
  this scale, they are generally smaller than the plotting markers.}
  \label{fig:storage-feedback-app}
\end{figure}

\FloatBarrier

\section{Supporting Mechanism and Frontier Results}
\label{app:mechanism-tables}

\subsection{Scalar-temperature control}

For positive categorical laws $\nu$ and $\mu$, define
$\nu_\beta(a)=\nu(a)^\beta/\sum_b \nu(b)^\beta$. Then
$\mu=\nu_\beta$ for some $\beta\ge0$ exactly when
\begin{equation}
 \log\frac{\mu(a)}{\mu(b)}=\beta\log\frac{\nu(a)}{\nu(b)}
 \qquad\text{for every }a,b.
 \label{eq:temperature-characterization-app}
\end{equation}
The forward direction is immediate. For the reverse direction, fix one
reference label $b$, exponentiate the equality, and normalize. Scalar
temperature preserves rankings and ties for $\beta>0$; the boundary
$\beta=0$ instead maps every positive $\nu$ to the uniform law and collapses all
rankings to ties. Define the irreducible temperature residual by
\begin{equation}
 R_{\mathrm{temp}}(\mu;\nu)
 :=\min_{\beta\ge0}\KL(\mu\kldiv \nu_\beta).
 \label{eq:temperature-residual-app}
\end{equation}
The objective
\[
 f(\beta)=-\Ent(\mu)-\beta\E_\mu[\log \nu(A)]
 +\log\sum_a \nu(a)^\beta
\]
for $A\sim\mu$ is convex because
$f''(\beta)=\operatorname{Var}_{\nu_\beta}[\log \nu(A)]\ge0$. The following
same-state audit minimizes this objective separately at every coordinate, a
stronger control than fitting one global temperature. None of the primary
fitted optima occurs at $\beta=0$.

\subsection{Same-state temperature and token values}

The same-state audit follows full-support \CRS{} ($H=\infty$) paths at 64 NFE. In
each of four shards it uses eight paths, seven post-warmup steps, and eight
deterministically spaced active coordinates per step, for 1792 primary
coordinate--step pairs. At pair $j$, let $h_j^{\mathrm{ctx}}$ be the native
categorical denoiser output under the current input overwrite and let $h_j^0$ be the
output from the identical mutable state and time after removing all current
input overwrites. The audited coordinate itself is uncommitted and
has the same input token in both inputs. We fit
\begin{equation}
 \widehat\beta_j
 :=\arg\min_{\beta\in[0,4096]}
 \KL\!\left(h_j^{\mathrm{ctx}}\kldiv(h_j^0)_\beta\right),
 \qquad
 R_j:=\KL\!\left(h_j^{\mathrm{ctx}}
 \kldiv(h_j^0)_{\widehat\beta_j}\right).
 \label{eq:same-state-fit-app}
\end{equation}
For shard $s$, the reported residual fraction is the pooled ratio
\begin{equation}
 \mathcal R_s
 :=\frac{\sum_{j\in\mathcal J_s}R_j}
 {\sum_{j\in\mathcal J_s}
 \KL(h_j^{\mathrm{ctx}}\kldiv h_j^0)}.
 \label{eq:same-state-pooled-residual-app}
\end{equation}
Float64 log probabilities and 48 bisection steps solve the convex fit; no
primary optimum reaches either boundary; in particular, none uses
$\beta=0$, the uniform-reference limit.
The mean and SE of $\mathcal R_s$ over the four shards are
$79.4\pm2.1\%$. The top-ranked token differs between
$h_j^{\mathrm{ctx}}$ and $h_j^0$ on $30.2\pm0.8\%$ of coordinates, and
$99.89\pm0.06\%$ contain a pairwise rank inversion in the union of their two
top-64 sets (log-probability tie tolerance $10^{-12}$). The initial-step
empty-context control has zero top-rank changes and zero rank inversions.

At full clean-token support, the value intervention holds the selected positions,
transition streams, and argmax decoding in the final step fixed. Argmax values give
$28.27\pm0.53$ GenPPL. Replacing them by the current second-ranked token gives
$77.75\pm1.80$, and sampling a top-$p=0.9$ alternative after excluding the
argmax gives $78.36\pm1.56$. The refreshed-value branch, which keeps positions
visible but updates values from each current argmax, gives
$26.42\pm0.39$. Token identity matters, while immutable identity is not
required for hard visible context to alter the path.

\subsection{Hidden-coordinate support}

The later-support audit uses four shards, eight \CRS{} ($H=\infty$) paths per
shard, and up to eight spaced commitment events per path. At offsets $1,2,4,$
and $8$, it evaluates the stored label under the overwritten input, the same
\CRS{} trajectory with only the audited coordinate restored to its mutable
noisy state, and a paired fully revisable trajectory. The restored-coordinate
view is an input intervention, not exact marginalization over that coordinate.

For the restored-coordinate view, stored-label probabilities at offsets
$1,2,4,8$ are $0.7483\pm0.0060$, $0.8323\pm0.0044$, $0.8230\pm0.0052$, and
$0.8291\pm0.0009$; paired revisable values are $0.4939\pm0.0262$,
$0.5542\pm0.0359$, $0.5481\pm0.0422$, and $0.5630\pm0.0430$. Restoration changes the
coordinate in $144,114,108,96$ of $256,224,224,223$ events. Means and standard
errors use four shards; direct recomputation from the shard files gives
$0.00089$ for the last restored-coordinate SE. Conditional on the restoration
actually changing that coordinate, the restored probabilities are
$0.7300\pm0.0129$, $0.8402\pm0.0209$, $0.8199\pm0.0297$, and
$0.8458\pm0.0230$; the paired revisable probabilities are
$0.4534\pm0.0327$, $0.5295\pm0.0493$, $0.5191\pm0.0486$, and
$0.5378\pm0.0580$. Because other visible context and noisy-state changes
remain, the audit supports a trajectory-mediated rather than isolated account.

\iffalse
\subsection{Adaptive current-step operators}

Table~\ref{tab:operator-robustness-app} reports a completed 256-sample CUDA
screen. The three operators recompute within-step support from each arm's own
trajectory at every post-warmup step and all active coordinates. Let
$p_{(1)}\ge p_{(2)}$ be the two largest denoiser-output probabilities. The absolute
margin arm uses
\begin{equation}
 p_{r,i}=1-(1-0.8)s_{r,i},
 \qquad s_{r,i}=p_{(1)}-p_{(2)},
 \label{eq:adaptive-margin-app}
\end{equation}
so $p_{r,i}\in[0.8,1.0]$. The binary-information arm uses the same map with
\begin{equation}
 s_{r,i}=(p_{(1)}+p_{(2)})
 \left[1-\frac{H_2\!\left(p_{(1)}/(p_{(1)}+p_{(2)})\right)}{\log 2}\right],
 \quad H_2(u)=-u\log u-(1-u)\log(1-u).
 \label{eq:adaptive-binary-app}
\end{equation}
The charge-capped arm computes the converted-clean charge of global $p=0.8$ at
the same state, allocates that charge across coordinates in proportion to
absolute margin, and chooses the strongest prefix that does not exceed each
allocation. It obeys its upper budget but uses less than ten percent of the
global-$p=0.8$ reference charge, so it is an audit rather than a matched-charge
baseline. Each row compares the $H=8$ and $H=0$ policies for that operator.

\begin{table}[H]
\centering
\small
\setlength{\tabcolsep}{5pt}
\begin{tabular}{@{}lrrr@{}}
\toprule
Current-step operator & $\Delta$ GPT-2 PPL & $\Delta$ Pythia PPL & $\Delta$ Rep-4 \\
\midrule
Absolute margin & $-6.98\pm0.14$ & $-6.39\pm0.11$ & $+0.00691\pm0.00083$ \\
Binary information & $-7.89\pm0.19$ & $-7.15\pm0.10$ & $+0.00770\pm0.00070$ \\
Charge-capped absolute & $-6.43\pm0.33$ & $-5.96\pm0.28$ & $+0.00533\pm0.00050$ \\
\bottomrule
\end{tabular}
\caption{Eight-step minus zero-step changes across adaptive current-step
operators. Negative PPL and positive Rep-4 show the recurring likelihood and
concentration movement.}
\label{tab:operator-robustness-app}
\end{table}
\fi

\subsection{Matched GenPPL--entropy comparisons}

On Duo-distilled, the full-support \CRS{} ($H=\infty$) sweep varies the confidence budget and
compares each point with a global top-$p$ curve using the same argmax decoding in the final step,
$p\in\{0.75,0.8,0.85,0.9,0.95,1.0\}$. For each shard we interpolate that
shard's observed global curve at that shard's observed CRS diversity and form
the residual before taking the four-shard mean and SE; no extrapolation is
used. Both sweeps were generated in the same launch with aligned pair seeds,
initial states, and RNG resets across cells, so the four interpolated values
are paired shard contrasts. The entropy-matched GenNLL residual retains its
sign across the budget sweep: for
$D_{\max}=128,256,512,1024$, respectively, it is
$-0.089\pm0.029$, $-0.130\pm0.038$, $-0.154\pm0.031$, and
$-0.155\pm0.030$ nats per token.
At $D_{\max}=512$, matching separately on each diversity statistic gives the following table,
whose entries are CRS-minus-global-top-$p$ GenNLL residuals in nats per
token; each column names the diversity statistic used for matching. Negative
values favor CRS.
\begin{center}
\begin{tabular}{@{}lrrrr@{}}
\toprule
Matched summary & Entropy & Dist-4 & Rep-4 & Self-BLEU \\
\midrule
GenNLL residual & $-0.154\pm0.031$ & $-0.025\pm0.048$ & $+0.070\pm0.036$ & $-0.144\pm0.019$ \\
\bottomrule
\end{tabular}
\end{center}


\iffalse
\section{Additional Support and Endpoint Diagnostics}
\label{app:spatial-diagnostics}

\paragraph{Clean-hypothesis coverage.}
We forward-corrupt held-out clean text to $t=0.8$ and inspect the same-state
top-$p$ support. At $p=0.8$, the top-$p$ support contains the true clean token at
$0.7173\pm0.0041$ of coordinates, excluding $289.4\pm4.2$ true tokens per
length-$1024$ sequence. The probability that an entire contiguous 16-token
clean block remains covered is $0.0214\pm0.0014$. This is a same-state
coverage diagnostic; it does not estimate generated-path reachability.

\paragraph{Enumerated-target toy.}
The point-codebook target has vocabulary size eight, sequence length five, and
eight deterministic codewords with probabilities
$(0.3092,\allowbreak 0.1776,\allowbreak 0.1284,\allowbreak 0.1020,\allowbreak
0.0853,\allowbreak 0.0737,\allowbreak 0.0652,\allowbreak 0.0586)$. Exact posterior marginals and coordinatewise uniform
bridges are recomputed over a 64-step cosine schedule, and native arms sample
every transition. Toy CRS stores two model-selected LOO argmax values after the
21-step warmup, overwrites later inputs, and copies them into the endpoint;
other coordinates retain their sampled terminal states. This differs from the
argmax decoding in the final step used in the learned-model factorial.

From 200,000 endpoints, we report unsmoothed plug-in
$\KL(q_0\kldiv\widehat P)$ and total variation (TV); all eight target states are observed. Native
full support gives $0.00010/0.0060$, top-$p=0.9$ and $0.8$ give
$0.00111/0.0198$ and $0.00433/0.0358$, and fixed two-commit CRS gives
$0.1248/0.1546$. The target is exact and endpoint divergences are Monte Carlo
estimates. The toy shows that fixed truncation and commitment can leave
endpoint bias; it is not a ranking of the learned text samplers.
\fi


\end{document}
